\documentclass{article}
\usepackage{algorithm}
\usepackage{algorithmic}
% if you need to pass options to natbib, use, e.g.:
%     \PassOptionsToPackage{numbers, compress}{natbib}
% before loading neurips_2026

% The authors should use one of these tracks.
% Before accepting by the NeurIPS conference, select one of the options below.
% 0. "default" for submission
\usepackage{neurips_2026}
% the "default" option is equal to the "main" option, which is used for the Main Track with double-blind reviewing.
% 1. "main" option is used for the Main Track
%  \usepackage[main]{neurips_2026}
% 2. "position" option is used for the Position Paper Track
%  \usepackage[position]{neurips_2026}
% 3. "eandd" option is used for the Evaluations & Datasets Track
 % \usepackage[eandd]{neurips_2026}
 % if you need to opt-in for a single-blind submission in the E&D track:
 %\usepackage[eandd, nonanonymous]{neurips_2026}
% 4. "creativeai" option is used for the Creative AI Track
%  \usepackage[creativeai]{neurips_2026}
% 5. "sglblindworkshop" option is used for the Workshop with single-blind reviewing
 % \usepackage[sglblindworkshop]{neurips_2026}
% 6. "dblblindworkshop" option is used for the Workshop with double-blind reviewing
%  \usepackage[dblblindworkshop]{neurips_2026}

% After being accepted, the authors should add "final" behind the track to compile a camera-ready version.
% 1. Main Track
 % \usepackage[main, final]{neurips_2026}
% 2. Position Paper Track
%  \usepackage[position, final]{neurips_2026}
% 3. Evaluations & Datasets Track
 % \usepackage[eandd, final]{neurips_2026}
% 4. Creative AI Track
%  \usepackage[creativeai, final]{neurips_2026}
% 5. Workshop with single-blind reviewing
%  \usepackage[sglblindworkshop, final]{neurips_2026}
% 6. Workshop with double-blind reviewing
%  \usepackage[dblblindworkshop, final]{neurips_2026}
% Note. For the workshop paper template, both \title{} and \workshoptitle{} are required, with the former indicating the paper title shown in the title and the latter indicating the workshop title displayed in the footnote.
% For workshops (5., 6.), the authors should add the name of the workshop, "\workshoptitle" command is used to set the workshop title.
% \workshoptitle{WORKSHOP TITLE}

% "preprint" option is used for arXiv or other preprint submissions
 % \usepackage[preprint]{neurips_2026}

% to avoid loading the natbib package, add option nonatbib:
%    \usepackage[nonatbib]{neurips_2026}

\usepackage[utf8]{inputenc} % allow utf-8 input
\usepackage[T1]{fontenc}    % use 8-bit T1 fonts
\usepackage{hyperref}       % hyperlinks
\usepackage{url}            % simple URL typesetting
\usepackage{booktabs}       % professional-quality tables
\usepackage{amsfonts}       % blackboard math symbols
\usepackage{nicefrac}       % compact symbols for 1/2, etc.
\usepackage{microtype}      % microtypography
\usepackage{xcolor}         % colors

%add package

\usepackage{multirow}
\usepackage{colortbl}
\usepackage{subfigure}
\usepackage{amsmath}
\usepackage{amssymb}
\usepackage{wrapfig}
\usepackage{enumitem}
\usepackage{amsthm}
\usepackage{graphicx}
\usepackage{array}
\usepackage{framed}
\definecolor{promptbg}{RGB}{246,246,250}
\definecolor{prompttitle}{RGB}{40,70,135}
\colorlet{shadecolor}{promptbg}

% Note. For the workshop paper template, both \title{} and \workshoptitle{} are required, with the former indicating the paper title shown in the title and the latter indicating the workshop title displayed in the footnote. 
\title{KnowledgeClaw: Towards Agentic Knowledge Engineering with Omnimodal Hypergraphs}


% The \author macro works with any number of authors. There are two commands
% used to separate the names and addresses of multiple authors: \And and \AND.
%
% Using \And between authors leaves it to LaTeX to determine where to break the
% lines. Using \AND forces a line break at that point. So, if LaTeX puts 3 of 4
% authors names on the first line, and the last on the second line, try using
% \AND instead of \And before the third author name.


\author{%
  David S.~Hippocampus\thanks{Use footnote for providing further information
    about author (webpage, alternative address)-\emph{not} for acknowledging
    funding agencies.} \\
  Department of Computer Science\\
  Cranberry-Lemon University\\
  Pittsburgh, PA 15213 \\
  \texttt{hippo@cs.cranberry-lemon.edu} \\
  % examples of more authors
  % \And
  % Coauthor \\
  % Affiliation \\
  % Address \\
  % \texttt{email} \\
  % \AND
  % Coauthor \\
  % Affiliation \\
  % Address \\
  % \texttt{email} \\
  % \And
  % Coauthor \\
  % Affiliation \\
  % Address \\
  % \texttt{email} \\
  % \And
  % Coauthor \\
  % Affiliation \\
  % Address \\
  % \texttt{email} \\
}


\begin{document}


\maketitle

\vspace{-4mm}


\begin{abstract}
Knowledge engineering (KE) has progressed from chunk-based Retrieval-Augmented Generation (RAG) to graph-based RAG, and more recently to agentic systems capable of multimodal ingestion and lifecycle maintenance. However, existing agentic KE systems remain incomplete: each satisfies at most two of three necessary properties: omnimodality, closed-loop lifecycle, and multi-angle application, limiting their ability to serve a live, all-modality knowledge base. In this work, we propose \textbf{KnowledgeClaw}, the first agentic knowledge engineering system that jointly satisfies all three properties on a unified omnimodal knowledge hypergraph. KnowledgeClaw consists of a comprehensive function pipeline, including omnimodal knowledge hypergraph construction, closed-loop lifecycle maintenance, and multi-angle application with intent routing. Experiments on our released OmniKE-Bench against 22 baselines show that KnowledgeClaw obtains the highest score on all six OmniKE tasks, attaining an aggregate OKE of 59.8 and surpassing the strongest baseline by +13.6. Our software and data are available.\footnote{Anonymous GitHub link: \url{https://anonymous.4open.science/r/KnowledgeClaw}}.
\end{abstract}

\section{Introduction}\label{sec:intro}

\begin{wrapfigure}{r}{0.54\textwidth}
\vspace{-4.5mm}
\centering
\includegraphics[width=0.54\textwidth]{F1.pdf}
\caption{An illustration of KnowledgeClaw agent.}\label{F1}
\vspace{-2mm}
\end{wrapfigure}

Knowledge engineering (KE)~\citep{KEclassic} aims to transform unstructured information into queryable, reasoning-ready system, but it relies on manual expert construction for a long time and resists large-scale deployment~\citep{KGLLMsurvey}. Large Language Models (LLMs)~\citep{GPT4,LLMsurvey} and autonomous agents~\citep{ReAct,AgentSurvey} loosen this constraint: an agent can now read files, extract entities, link relations, reason, and update knowledge end to end~\citep{LLMIEsurvey,KGLLMsurvey}, making it feasible to continuously curate raw files into a live knowledge base.

Realizing agentic knowledge engineering in practice is demanding. Consider a researcher's knowledge base: papers arrive as PDFs, lectures as audio, demos as video, slides as images, and notes as text; retractions and revisions rewrite earlier entries; and queries range from locating a sentence in a recorded talk to surfacing unexplored directions. Such workloads demand three properties: \textbf{(i) Omnimodality}: a unified index over text, file, image, audio, and video~\citep{Qwen3-Omni,Qwen3.5-Omni,OmniRAG-Agent}; \textbf{(ii) Closed-loop lifecycle}: autonomous consistency under add/delete/refresh operations~\citep{MemorySurvey,MemGPT,Mem0,A-MEM}; and \textbf{(iii) Multi-angle application}: joint support for precise lookup, topic synthesis, and proactive discovery~\citep{GraphRAG,HyperGraphRAG,LLMWiki}.

\begin{figure}[t]
\centering
\includegraphics[width=13.8cm]{F2.pdf}
\caption{Three contributions map onto three necessary properties: C1 unifies five modalities via cross-modal hyperedges (T1, T2), C2 maintains lifecycle consistency on the live hypergraph (T5), and C3 routes queries among subgraph retrieval, community synthesis, and knowledge discovery (T3, T4, T6). KnowledgeClaw obtains the highest scoring on all six OmniKE tasks (OKE\,$=$\,59.8).}\label{F2}
\vspace{0.5mm}
\end{figure}

However, existing systems each rarely satisfy all of these three properties. Chunk- and graph-based RAG~\citep{RAG,RAPTOR,GraphRAG,LightRAG,HippoRAG2,HyperGraphRAG}, ontology-grounded knowledge graph (KG) construction~\citep{iText2KG,OntoKG}, and agentic KG engines~\citep{KAG,Graphiti} most operate on static text alone, violating (i) and (ii). Memory agents (MemGPT~\citep{MemGPT}, Mem0~\citep{Mem0}, A-MEM~\citep{A-MEM}) provide full CRUD but rely on flat vector storage without structured omnimodal representation, violating (i) and (iii). Omnimodal pipelines such as OmniRAG-Agent~\citep{OmniRAG-Agent} extend retrieval across modalities but lack lifecycle maintenance and multi-angle application, violating (ii) and (iii). OpenClaw~\citep{OpenClaw} and its successors Cognee~\citep{Cognee}, LLM Wiki~\citep{LLMWiki}, and Graphify~\citep{Graphify} approach unified agentic KE but still lack a common benchmark and evaluation that jointly measure all three properties, motivating the unified agent illustrated in Figure~\ref{F1}.

To address these gaps, for the first time, we formalize the problem as \emph{Omnimodal Knowledge Engineering} (OmniKE) and release \emph{OmniKE-Bench}, with six method-agnostic tasks (T1 knowledge acquisition, T2 cross-modal coreference, T3 multi-hop retrieval, T4 cross-modal reasoning, T5 lifecycle consistency, T6 knowledge discovery) that jointly stress the three axes across 10 professional-domain personas, 5 modalities, and $\approx$4{,}500 real files, with protocols adapted from previous traditional KE evaluation methods~\citep{RAGAS}. Figure~\ref{F2} previews how our contributions cover these tasks.

Building on this foundation, we propose \textbf{KnowledgeClaw}, an agentic system that autonomously maintains an omnimodal knowledge hypergraph through three contributions. \textbf{(C1) Omnimodal Knowledge Hypergraph} unifies entities from all five modalities (text, PDF, image, audio, video) under cross-modal hyperedges~\citep{Hypergraph,HAHE,Text2NKG}, extending HyperGraphRAG~\citep{HyperGraphRAG} beyond text. \textbf{(C2) Closed-loop Lifecycle Agent} runs a build-query-reason-update loop with consistency-preserving addition, cascading deletion, and orphan cleanup, lifting the static-corpus assumption~\citep{GraphRAG,LightRAG,HippoRAG2}. \textbf{(C3) Multi-angle Application with Intent Routing} fuses subgraph-level retrieval~\citep{HyperGraphRAG,PathRAG,HippoRAG2}, community-level synthesis via Louvain~\citep{Louvain} clustering with summary embeddings~\citep{GraphRAG}, and community-based knowledge discovery under a lightweight intent router~\citep{SelfRAG,AdaptiveRAG}, each targeting a distinct query type.

Experiments on OmniKE-Bench against 22 baselines spanning seven paradigm groups show that KnowledgeClaw is the only system to score on all six tasks and ranks first on each (T1 = 74.9, T2 = 89.3, T3 = 34.5, T4 = 48.9, T5 = 46.8, T6 = 64.3), attaining an OKE of 59.8 that surpasses the strongest external baseline LLM Wiki~\citep{LLMWiki} by \textbf{+13.6} and our prior anchor OpenClaw~\citep{OpenClaw} by \textbf{+15.4}. Three orthogonal ablations confirm each contribution precisely targets the property it serves, demonstrating KnowledgeClaw's potential for real-world agentic knowledge engineering.

\clearpage

\begin{figure}[t]
\centering
\includegraphics[width=\textwidth]{F3.pdf}
\caption{\textbf{(A)} Six OmniKE tasks covering construction (T1 acquisition, T2 cross-modal coreference), maintenance (T5 lifecycle consistency), and application (T3 multi-hop, T4 cross-modal, T6 discovery). \textbf{(B)} OmniKE-Bench: 10 personas, 10 domains, 5 modalities, and 4{,}544 real files.}\label{F3}
\end{figure}

\section{The OmniKE Tasks and Benchmark}\label{sec:task}

We formalize \emph{Omnimodal Knowledge Engineering} (\textbf{OmniKE}) as six method-agnostic tasks (Figure~\ref{F3}A) that jointly cover the three properties \textbf{(i)--(iii)} of \S\ref{sec:intro}, and instantiate them as \emph{OmniKE-Bench} (Figure~\ref{F3}B), the first end-to-end benchmark for agentic KE on real omnimodal personal files.

\textbf{Task Definition.} The six tasks form three groups, one per property. \emph{Construction (T1, T2)} targets \textbf{(i) Omnimodality} over $\mathcal{M}\!=\!\{\textit{text},\textit{pdf},\textit{image},\textit{audio},\textit{video}\}$: given files $\mathcal{D}\!=\!\{f_i\}$ with $m(f_i)\!\in\!\mathcal{M}$, \emph{T1 (Knowledge Acquisition)} outputs entities $\hat{\mathcal{E}}$, propositions $\hat{\mathcal{P}}$, and modality map $\hat{\mu}\!:\!\hat{\mathcal{E}}\!\to\!2^{\mathcal{M}}$, and \emph{T2 (Cross-modal Coreference)} links mentions $(e_a,e_b)$ to a canonical referent across modalities. \emph{Maintenance (T5)} targets \textbf{(ii) Closed-loop Lifecycle} via \emph{NL CRUD} (e.g., ``this paper is retracted, remove it'' + verification probe) and \emph{Fact-refresh} (old fact, its replacement, and a question whose answer must track the change). \emph{Application (T3, T4, T6)} targets \textbf{(iii) Multi-angle Application}: \emph{T3 (Multi-hop Retrieval)} returns an answer $\hat{a}$ with evidence for a 2--4-hop question $q$; \emph{T4 (Cross-modal Reasoning)} restricts T3 to answers \emph{requiring} evidence from $\geq\!2$ modalities; \emph{T6 (Knowledge Discovery)} emits \emph{known}($t$), \emph{gaps}($t$), and \emph{sug}($t$) for a topic $t$. The aggregate \emph{OKE Score} averages T1--T6 on $0$--$100$, scoring architecturally unsupported tasks as $0$ so partial coverage of (i)--(iii) is penalized.

\textbf{Benchmark Construction.} OmniKE-Bench is organized by \emph{personas}, each a researcher with months of accumulated files in one professional domain, spanning \textbf{10} domains (Healthcare, Finance, Education, Law, Psychology, Music, Gaming, Building, Software, Robotics), \textbf{5} modalities, and \textbf{4{,}544} real native files ($\approx$22\,GB; 985\,/\,1{,}285\,/\,1{,}507\,/\,362\,/\,255 for text\,/\,PDF\,/\,image\,/\,audio\,/\,video). Construction runs in four phases (Figure~\ref{F3}B): \emph{(1) Corpus}: collect real files per persona and modality; \emph{(2) Derive}: convert each file to a unified JSON record via PyMuPDF (PDFs), \texttt{gpt-4o-mini} vision (image captions), and \texttt{gemini-3.1-flash-lite} transcription (audio/video after \texttt{ffmpeg} downsampling); \emph{(3) Annotate}: generate gold with \texttt{gpt-4o-mini} for T3/T4/T5 (NL-CRUD + fact-refresh probes) and T6 ($\sim$35 topics/persona), leaving T1, T2 reference-free; \emph{(4) Freeze}: package each persona for release. Total cost: $\approx$30~min derivation + $\approx$\$30 annotation, $\sim$60\% cheaper than annotating all six tasks.

\textbf{Evaluation Method.} Adapting traditional KE evaluation methods~\citep{RAGAS}, all tasks score \emph{automatically} without LLM calls ($<$1\,s/run); reference-free tasks (T1, T2, T6) score against the corpus, not gold labels. \emph{T1}: entity/proposition faithfulness (substring or $>\!20\%$ token overlap), quantity sanity, and modality coverage. \emph{T2}: link faithfulness, semantic relatedness, non-triviality, modality diversity, and link quantity. \emph{T3}: SQuAD-style token F1. \emph{T4}: $0.7\!\cdot\!\text{F1}\!+\!0.3\!\cdot\!\text{XMEv}$ (cross-modal evidence rate). \emph{T5}: NL-CRUD accuracy (keyword check; delete\,/\,update\,/\,others $3\!:\!2\!:\!1$) + fact-refresh F1. \emph{T6}: $\geq\!3$-char bigram presence in KB (\emph{known} should appear, \emph{gaps} should not).

\begin{figure}[t]
\centering
\includegraphics[width=\textwidth]{F4.pdf}
\caption{\textbf{KnowledgeClaw} methodology, where C1--C3 build, maintain, and serve an omnimodal knowledge hypergraph. \textbf{(A)} C1 extracts entities from text, PDF, image, audio, and video and fuses them via within-modality canonicalization and cross-modal coreference into typed hyperedges. \textbf{(B)} C2 closes the lifecycle loop via a Build--Query--Reason--Update pass with orphan eviction and fact refresh. \textbf{(C)} C3 routes queries among subgraph retrieval, community synthesis, and discovery.}\label{F4}
\end{figure}

\section{Methodology: KnowledgeClaw}\label{sec:method}

\textbf{KnowledgeClaw} (Figure~\ref{F4}) realizes \textbf{(i)--(iii)} via three contributions that respectively \emph{build} (C1), \emph{maintain} (C2), and \emph{serve} (C3) an omnimodal knowledge hypergraph; notation follows \S\ref{sec:task}.

\subsection{C1: Omnimodal Knowledge Hypergraph}\label{sec:c1}

To unify entities from all five modalities under a single algebraic structure, C1 lifts the substrate from a binary graph to a typed hypergraph:
\begin{equation}
\mathcal{G}\!=\!\bigl(\mathcal{E},\mathcal{H},\mathcal{A},\mathcal{P};\,\mu,\tau,\chi,\phi\bigr),\quad
h\!=\!\bigl(\mathcal{E}_h,\tau_h,\pi_h,\mathcal{A}_h,\omega_h\bigr),\;\,|\mathcal{E}_h|\!\geq\!2,\;\,\omega_h\!\in\![0,1],
\label{eq:hg}
\end{equation}
where $\mathcal{E}$ is the entity set, $\mathcal{H}\!\subseteq\!2^{\mathcal{E}}$ the typed hyperedges, $\mathcal{A}\!\subseteq\!\mathcal{D}$ the resident assets, and $\mathcal{P}$ the propositions; the four maps $\mu\!:\!\mathcal{E}\!\to\!2^{\mathcal{M}}$, $\tau\!:\!\mathcal{H}\!\to\!\mathcal{T}$, $\chi\!:\!\mathcal{E}^{2}\!\to\!\{0,1\}$, $\phi\!:\!\mathcal{E}\!\cup\!\mathcal{H}\!\to\!\mathbb{R}^{d}$ assign each entity its observed modalities (extending $\hat{\mu}$), each hyperedge a relation type, each entity pair a cross-modal coreference flag, and every node and hyperedge a shared embedding; each $h$ further records its proposition $\pi_h$, supporting assets $\mathcal{A}_h$, and confidence $\omega_h$. Allowing $|\mathcal{E}_h|\!\geq\!2$ with mixed-modality $\mathcal{A}_h$, Eq.~\eqref{eq:hg} subsumes binary graph RAG~\citep{GraphRAG,LightRAG} as the $|\mathcal{E}_h|\!=\!2$ case and lifts HyperGraphRAG~\citep{HyperGraphRAG} from text to all five modalities.

Populating $\mathcal{G}$ proceeds in two stages. An LLM extractor first emits $(\hat{\mathcal{E}}_f,\hat{\mathcal{P}}_f,\hat{\mathcal{H}}_f,\hat{\mu}_f)\!=\!\mathrm{Extract}_{\mathrm{LLM}}(f,m(f))$ for every $f\!\in\!\mathcal{A}$, and a canonicalization $\kappa\!:\!\bigsqcup_f\!\hat{\mathcal{E}}_f\!\to\!\mathcal{E}$ merges surface variants within each modality. Mentions that $\kappa$ cannot bridge across modalities are linked by an explicit coreference test, defined for any $e_a,e_b$ with $\mu(e_a)\!\cap\!\mu(e_b)\!=\!\emptyset$ as:
\begin{equation}
\chi(e_a,e_b)\!=\!\mathbb{1}\!\Bigl[\,\bigl\langle\phi(e_a),\phi(e_b)\bigr\rangle\!\geq\!\theta_{\chi}\;\wedge\;\mathrm{Verify}_{\mathrm{LLM}}\!\bigl(e_a,e_b\,\big|\,\mathcal{S}(e_a)\!\cup\!\mathcal{S}(e_b)\bigr)\,\Bigr],
\label{eq:coref}
\end{equation}
where $\langle\cdot,\cdot\rangle$ is cosine similarity, $\theta_{\chi}$ a tunable margin, and $\mathcal{S}(e)\!=\!\bigcup_{h:\,e\in\mathcal{E}_h}\!\mathcal{A}_h$ the source-asset support; the verifier rejects pairs that pass cosine but fail semantic agreement. Together, $\kappa$ fuses within-modality mentions while $\chi$ links across modalities, jointly grounding T1 and T2.

\textbf{Proposition 1.} \textit{$\mathcal{G}$ strictly subsumes text-only graphs: $n$-ary cross-modal facts on any $\mathcal{M}'\!\subseteq\!\mathcal{M}$ are uniquely representable in $\mathcal{G}$ yet unrecoverable from text, generalizing uniformly across modalities.}\vspace{-2mm}
\begin{proof}
We provide experimental results in Sections~\ref{5.4} and proofs in Appendix~\ref{proof1}.
\end{proof}

\subsection{C2: Closed-loop Lifecycle Agent}\label{sec:c2}

A live knowledge base stays correct only if every corpus change propagates to $\mathcal{G}$. Each event $\Delta_n\!\in\!\{\textsc{Add},\textsc{Del},\textsc{Refresh}\}\!\times\!\mathcal{D}$ at step $n$ drives a four-stage \emph{Build--Query--Reason--Update} pass:
\begin{equation}
\mathcal{G}_{n+1}\!=\!\bigl(\mathrm{Update}\!\circ\!\mathrm{Reason}\!\circ\!\mathrm{Query}\!\circ\!\mathrm{Build}\bigr)\!\bigl(\mathcal{G}_n,\Delta_n\bigr),\;\;
\mathcal{L}_{n+1}\!=\!\mathcal{L}_n\!\cup\!\bigl\{\bigl(\Delta_n,\mathcal{G}_{n+1}\!\ominus\!\mathcal{G}_n\bigr)\bigr\},
\label{eq:loop}
\end{equation}
where the four stages respectively ingest $\Delta_n$, retrieve its neighborhood, resolve conflicts via LLM, and commit $\mathcal{G}_{n+1}$; $\ominus$ is the structural diff between consecutive states and $\mathcal{L}_n$ a replayable audit log. The non-trivial operator $\textsc{Del}$ avoids both dangling entities and over-deletion: reusing $\mathcal{S}(\cdot)$ from Eq.~\eqref{eq:coref}, we evict orphans with their unsupported hyperedges in one transactional pass:
\begin{equation}
\mathcal{E}_n^{\mathrm{orph}}(f)\!=\!\bigl\{\,e\!\in\!\mathcal{E}_n\!:\,\mathcal{S}(e)\!\subseteq\!\{f\}\,\bigr\},\;\;
\begin{aligned}[t]
\mathcal{E}_{n+1} &= \mathcal{E}_n\!\setminus\!\mathcal{E}_n^{\mathrm{orph}}(f),\\[-1pt]
\mathcal{H}_{n+1} &= \bigl\{h\!\in\!\mathcal{H}_n\!:\,\mathcal{A}_h\!\not\subseteq\!\{f\}\,\wedge\,\mathcal{E}_h\!\cap\!\mathcal{E}_n^{\mathrm{orph}}(f)\!=\!\emptyset\bigr\},\\[-1pt]
\mathcal{A}_{n+1} &= \mathcal{A}_n\!\setminus\!\{f\},
\end{aligned}
\label{eq:cascade}
\end{equation}
where $\mathcal{E}_n^{\mathrm{orph}}(f)$ are entities supported solely by $f$, and the three updates drop these orphans, dangling hyperedges, and $f$, enforcing the cascading-delete invariant. $\textsc{Refresh}$ is the dual: it rewrites $h_{\mathrm{old}}$ to $h_{\mathrm{new}}$ with new proposition $\pi_{h_{\mathrm{new}}}$ and re-embeds in place,
\begin{equation}
\mathcal{H}_{n+1}\!=\!\bigl(\mathcal{H}_n\!\setminus\!\{h_{\mathrm{old}}\}\bigr)\!\cup\!\{h_{\mathrm{new}}\},\;\;
\phi(h_{\mathrm{new}})\!=\!\mathrm{Embed}\!\bigl(\pi_{h_{\mathrm{new}}}\bigr),
\label{eq:refresh}
\end{equation}
where $\mathrm{Embed}(\cdot)$ shares the encoder of $\phi$, so the neighborhood is preserved while retrieval tracks the new fact. Eqs.~\eqref{eq:loop}--\eqref{eq:refresh} lift the static-corpus assumption~\citep{GraphRAG,LightRAG,HippoRAG2,KAG} and ground T5.

\textbf{Proposition 2.} \textit{Under any finite event sequence, the lifecycle agent preserves consistency: $\mathcal{S}(e)\!\neq\!\emptyset$ for every $e\!\in\!\mathcal{E}_N$ and $\mathcal{E}_h\!\subseteq\!\mathcal{E}_N$ for every $h\!\in\!\mathcal{H}_N$.}\vspace{-2mm}
\begin{proof}
We provide experimental results in Sections~\ref{5.5} and proofs in Appendix~\ref{proof2}.
\end{proof}

\subsection{C3: Multi-angle Application with Intent Routing}\label{sec:c3}

KnowledgeClaw exposes three angles $\mathcal{R}\!=\!\{\mathrm{sub},\mathrm{com},\mathrm{disc}\}$ and routes among them by intent.

\paragraph{Subgraph-level retrieval ($\mathrm{sub}$).} For factoid/multi-hop $q$, dual indexes over $\mathcal{E},\mathcal{H}$ are fused via reciprocal rank fusion and top-$K$ seeds $\mathcal{Z}_q$ expanded by an $\ell$-hop walk~\citep{PathRAG,HippoRAG2}:
\begin{equation}
s_{\mathrm{RRF}}(x;q)\!=\!\!\!\!\sum_{r\,\in\,\{r_{\mathcal{E}},r_{\mathcal{H}}\}}\!\!\!\frac{1}{k_0\!+\!r(x;q)},\;\;\,
\mathcal{G}_q^{\mathrm{sub}}\!=\!\mathcal{G}\!\bigl[\mathcal{N}_{\leq\ell}(\mathcal{Z}_q)\bigr],\;\;\,
\mathcal{N}_{\leq\ell}(\mathcal{Z}_q)\!=\!\!\!\bigcup_{x\,\in\,\mathcal{Z}_q}\!\bigcup_{i=0}^{\ell}\!\mathcal{N}^{(i)}(x),
\label{eq:sub}
\end{equation}
with per-index ranks $r_{\mathcal{E}},r_{\mathcal{H}}$ under cosine to $\phi(q)$, RRF constant $k_0\!=\!60$, $i$-hop neighborhood $\mathcal{N}^{(i)}(x)$, and induced subgraph $\mathcal{G}[\,\cdot\,]$ (T3, T4).

\paragraph{Community-level synthesis ($\mathrm{com}$).} For topic queries, we partition the cooccurrence graph by Louvain modularity $Q$~\citep{Louvain}, summarize each $\mathcal{C}_k$ into $\sigma_k$ via LLM, and retrieve top-$L$:
\begin{equation}
Q\!=\!\frac{1}{2W}\!\!\sum_{i,j\,\in\,\mathcal{E}}\!\!\Bigl[A_{ij}\!-\!\tfrac{k_i k_j}{2W}\Bigr]\delta(c_i,c_j),\;\, A_{ij}\!=\!\!\!\sum_{h:\{i,j\}\subseteq\mathcal{E}_h}\!\!\!\!\omega_h,\;\,
\mathcal{C}_q^{\mathrm{com}}\!=\!\mathrm{TopL}_{\mathcal{C}_k}\!\bigl[\bigl\langle\phi(q),\phi(\sigma_k)\bigr\rangle\bigr],
\label{eq:com}
\end{equation}
with confidence-weighted adjacency $A_{ij}$ from $\mathcal{H}$, total weight $W\!=\!\tfrac{1}{2}\!\sum_{ij}\!A_{ij}$, weighted degree $k_i\!=\!\sum_j\!A_{ij}$, community label $c_i$, and Kronecker delta $\delta(\cdot,\cdot)$.

\paragraph{Knowledge discovery ($\mathrm{disc}$) and intent routing.} Reusing communities, for topic $t$ and query $q$ we emit the T6 triplet, score each angle, and concatenate selected outputs into $\hat{a}(q)$:
\begin{equation}
\begin{aligned}
\mathrm{known}(t)&\!=\!\bigl\{e\!\in\!\mathcal{E}\!:\,\langle\phi(t),\phi(e)\rangle\!\geq\!\theta_{\mathrm{k}}\bigr\},\;
\bigl(\mathrm{gaps}(t),\mathrm{sug}(t)\bigr)\!=\!\mathrm{LLM}_{\mathrm{disc}}\!\bigl(t,\{\sigma_k\}_{\mathcal{C}_t^{\mathrm{com}}},\mathrm{known}(t)\bigr),\\[-1pt]
\pi^{*}(q)&\!=\!\bigl\{a\!\in\!\mathcal{R}\!:\,w_a(q)\!\geq\!\theta_{\pi}\bigr\},\;\;
\hat{a}(q)\!=\!\mathrm{LLM}_{\mathrm{gen}}\!\Bigl(q,\!\!\bigoplus_{a\,\in\,\pi^{*}(q)}\!\!R_a(q)\Bigr),
\end{aligned}
\label{eq:disc-route}
\end{equation}
with discovery threshold $\theta_{\mathrm{k}}$, router score $w_a(q)\!=\!\mathrm{Router}(q,a)$ gated by $\theta_{\pi}$, concatenation $\bigoplus$, and per-angle outputs $R_{\mathrm{sub}}, R_{\mathrm{com}}, R_{\mathrm{disc}}$ from above, covering T3, T4, T6.

\textbf{Proposition 3.} \textit{Multi-angle routing $\pi^{*}(q)$ covers $\{T3,T4,T6\}$, while no single $a\!\in\!\mathcal{R}$ does.}\vspace{-2mm}
\begin{proof}
We provide experimental results in Sections~\ref{5.6} and proofs in Appendix~\ref{proof3}.
\end{proof}



\clearpage

\begin{table}[t]
\centering
\caption{Design comparison of all baselines across the six OmniKE tasks T1--T6: each cell summarizes the method's mechanism in text and formula, with ``-'' marking architecturally unsupported tasks. Rows follow Table~\ref{tab:main}'s paradigm grouping and order; notation follows \S\ref{sec:method}.}
\label{tab:design}
\setlength{\tabcolsep}{7.5pt}\renewcommand{\arraystretch}{1.2}
\resizebox{\textwidth}{!}{%
\begin{tabular}{l|cccccc}
\toprule
\textbf{Method} & \textbf{T1} & \textbf{T2} & \textbf{T3} & \textbf{T4} & \textbf{T5} & \textbf{T6} \\
\midrule
\rowcolor{cyan!5} NaiveRAG~\citep{RAG} & chunks $\{c_i\}$ & - & cosine top-$k$ & - & - & - \\
\rowcolor{cyan!5} RAPTOR~\citep{RAPTOR} & cluster tree & - & tree-walk top-$k$ & - & - & cluster $\sigma_k$ \\
\rowcolor{blue!7} GraphRAG~\citep{GraphRAG} & $G\!=\!(V,E)$ & - & comm.-routed top-$k$ & - & - & Louvain $\sigma_k$ \\
\rowcolor{blue!7} LightRAG~\citep{LightRAG} & dual-level $G$ & - & dual top-$k$ & - & - & - \\
\rowcolor{blue!7} HippoRAG2~\citep{HippoRAG2} & $G$ + PPR & - & PageRank top-$k$ & - & - & - \\
\rowcolor{blue!7} HyperGraphRAG~\citep{HyperGraphRAG} & $h\!\subseteq\!V$ & - & expand $\bigcup_h$ & - & - & - \\
\rowcolor{teal!8} iText2KG~\citep{iText2KG} & ontology $G$ & - & graph query & - & ontology rebuild & - \\
\rowcolor{teal!8} OntoKG~\citep{OntoKG} & ontology $G$ & - & SPARQL-like & - & - & - \\
\rowcolor{teal!14} KAG/OpenSPG~\citep{KAG} & live KG & - & KG-walk & - & partial CRUD & - \\
\rowcolor{teal!14} Graphiti~\citep{Graphiti} & temporal $G_t$ & - & temporal walk & - & partial CRUD & temporal $\sigma$ \\
\rowcolor{orange!8} MemGPT/Letta~\citep{MemGPT} & hier.\ mem $\{m_i\}$ & - & ANN & - & tool-call CRUD & - \\
\rowcolor{orange!8} Mem0~\citep{Mem0} & mem tree & - & hybrid top-$k$ & - & full CRUD & - \\
\rowcolor{orange!8} A-MEM~\citep{A-MEM} & agentic mem $m$ & - & structured & - & full CRUD & gap probe \\
\rowcolor{yellow!12} OmniRAG-Agent~\citep{OmniRAG-Agent} & $\{I_m\}_{m\in\mathcal{M}}$ & hashing & per-mod top-$k$ & mod.\ fusion & - & - \\
\rowcolor{green!8} OpenClaw~\citep{OpenClaw} & $G$+mem, omni & linking & top-$k$+mem & mod.\ concat & partial CRUD & LLM probe \\
\rowcolor{green!8} Cognee~\citep{Cognee} & $G$+mem, omni & linking & top-$k$+mem & mod.\ concat & partial CRUD & LLM probe \\
\rowcolor{green!8} LLM Wiki~\citep{LLMWiki} & $G$+mem, omni & linking & top-$k$+mem & mod.\ concat & partial CRUD & LLM probe \\
\rowcolor{green!8} Graphify~\citep{Graphify} & $G$+mem, omni & linking & top-$k$+mem & mod.\ concat & partial CRUD & LLM probe \\
\midrule
\rowcolor{red!8} KC w/o C1 & text $\mathrm{Extract}$ & - & RRF+$\ell$-hop & - & cascade & $\mathrm{LLM}_{\mathrm{disc}}$ \\
\rowcolor{red!8} KC w/o C2 & $\mathrm{Extract}\!+\!\kappa$ & $\chi$ test & RRF+$\ell$-hop & $\chi$-aware sub & - & $\mathrm{LLM}_{\mathrm{disc}}$ \\
\rowcolor{red!8} KC w/o C3 & $\mathrm{Extract}\!+\!\kappa$ & $\chi$ test & flat top-$k$ & flat top-$k$ & B-Q-R-U & flat top-$k$ \\
\rowcolor{red!12} \textbf{KnowledgeClaw} & $\mathrm{Extract}\!+\!\kappa$ & $\chi$ test & RRF+$\ell$-hop & $\chi$+$\bigoplus R_a$ & B-Q-R-U & $\mathrm{LLM}_{\mathrm{disc}}$ \\
\bottomrule
\end{tabular}
}
\end{table}

\section{Experiments}
\label{sec:exp}

We answer six RQs:
\textbf{RQ1:} Does KnowledgeClaw outperform existing methods on OmniKE-Bench?
\textbf{RQ2:} Does each contribution (C1, C2, C3) work?
\textbf{RQ3:} Does KnowledgeClaw generalize across the five modalities and ten personas?
\textbf{RQ4:} How robust is the lifecycle agent under continuous add/delete/refresh events?
\textbf{RQ5:} Does intent routing beat each individual retrieval mode on different queries?
\textbf{RQ6:} What are the time and cost of KnowledgeClaw in construction and serving?

\subsection{Experimental Setup}

\textbf{Datasets and Evaluation Metrics.} We evaluate on \emph{OmniKE-Bench} introduced in \S\ref{sec:task} (10 personas, 5 modalities, $\approx$4{,}500 real files); every task T1--T6 reports a 0--100 score and the aggregate \textbf{OKE} is their arithmetic mean, with architecturally unsupported tasks counted as 0 so partial coverage of properties (i)--(iii) is penalized. Per-persona scores average T1--T6 within a persona; full task definitions, data construction, and per-metric formulas are deferred to \S\ref{sec:task} and Appendix~E.

\textbf{Baselines and Ablations.} We compare KnowledgeClaw against 22 publicly available baselines spanning seven paradigm groups: \emph{chunk/hierarchical retrieval} (NaiveRAG~\citep{RAG}, RAPTOR~\citep{RAPTOR}); \emph{graph/hypergraph RAG} (GraphRAG~\citep{GraphRAG}, LightRAG~\citep{LightRAG}, HippoRAG2~\citep{HippoRAG2}, HyperGraphRAG~\citep{HyperGraphRAG}); \emph{KG construction} (iText2KG~\citep{iText2KG}, OntoKG~\citep{OntoKG}); \emph{agentic KG} (KAG/OpenSPG~\citep{KAG}, Graphiti~\citep{Graphiti}); \emph{memory agents} (MemGPT/Letta~\citep{MemGPT}, Mem0~\citep{Mem0}, A-MEM~\citep{A-MEM}); \emph{omnimodal RAG} (OmniRAG-Agent~\citep{OmniRAG-Agent}); and \emph{agentic KE} (OpenClaw~\citep{OpenClaw}, Cognee~\citep{Cognee}, LLM Wiki~\citep{LLMWiki}, Graphify~\citep{Graphify}). We additionally evaluate three internal ablations \emph{KC w/o C1/C2/C3}, each disabling exactly one contribution. All methods share the same input corpus, generation prompt, and LLM backbone for fairness; per-baseline adaptation details are in Appendix~E.

\textbf{Implementation Details.} We adopt \texttt{gemini-3.1-flash-lite-preview} as the unified LLM backbone and \texttt{gemini-embedding-2-preview} for all $\phi(\cdot)$. Subgraph retrieval uses RRF top-$K_{\mathcal{E}}\!=\!K_{\mathcal{H}}\!=\!40$ with walk depth $\ell\!=\!2$ and constant $k_0\!=\!60$; community synthesis runs Louvain at $\gamma\!=\!1.0$ with chunk size 300 tokens and threshold $\theta_{\chi}\!=\!0.78$; the intent router thresholds at $\theta_{\pi}\!=\!0.5$. Baselines run with default hyperparameters; all scores are averaged over five random seeds per persona on a single 80-core 512\,GB CPU server (no GPU; all LLM calls remote).

\begin{table}[t]
\centering
\caption{Main results on OmniKE-Bench. \textbf{T1}-\textbf{T6}: per-task scores; \textbf{HC}-\textbf{Rob}: per-persona averages over T1-T6; \textbf{OKE}: arithmetic mean over T1-T6 with architecturally unsupported tasks counted as 0; $\boldsymbol{\Delta}$: gap to OpenClaw (anchor); ``-'' = not architecturally supported. Best scores in bold.}
\label{tab:main}
\setlength{\tabcolsep}{3pt}\renewcommand{\arraystretch}{1.2} 
\resizebox{\textwidth}{!}{%
\begin{tabular}{l|cccccc|cccccccccc|cc}
\toprule
\textbf{Method} & \textbf{T1} & \textbf{T2} & \textbf{T3} & \textbf{T4} & \textbf{T5} & \textbf{T6} & \textbf{HC} & \textbf{Fin} & \textbf{Edu} & \textbf{Law} & \textbf{Psy} & \textbf{Mus} & \textbf{Gam} & \textbf{Bld} & \textbf{Sft} & \textbf{Rob} & \textbf{OKE} & $\boldsymbol{\Delta}$ \\
\midrule
\rowcolor{cyan!5} NaiveRAG~\citep{RAG} & 23.6 & - & 30.1 & - & - & - & 26.3 & 27.5 & 26.6 & 26.3 & 26.2 & 27.0 & 27.1 & 27.2 & 26.6 & 27.5 & 9.0 & \textcolor{red}{-35.4} \\
\rowcolor{cyan!5} RAPTOR~\citep{RAPTOR} & 26.7 & - & 33.3 & - & - & 51.5 & 33.6 & 35.8 & 37.8 & 41.7 & 40.3 & 37.9 & 34.8 & 41.5 & 32.9 & 35.8 & 18.6 & \textcolor{red}{-25.8} \\
\rowcolor{blue!7} GraphRAG~\citep{GraphRAG} & 34.4 & - & 34.1 & - & - & 58.2 & 40.4 & 40.7 & 44.5 & 45.9 & 44.7 & 42.3 & 39.8 & 45.5 & 38.4 & 40.5 & 21.1 & \textcolor{red}{-23.3} \\
\rowcolor{blue!7} LightRAG~\citep{LightRAG} & 32.8 & - & 33.8 & - & - & - & 32.8 & 32.9 & 33.7 & 32.5 & 32.8 & 33.9 & 33.8 & 33.7 & 32.9 & 34.1 & 11.1 & \textcolor{red}{-33.3} \\
\rowcolor{blue!7} HippoRAG2~\citep{HippoRAG2} & 37.6 & - & 34.4 & - & - & - & 35.6 & 34.9 & 36.7 & 35.0 & 35.3 & 36.8 & 36.4 & 36.4 & 36.0 & 36.9 & 12.0 & \textcolor{red}{-32.4} \\
\rowcolor{blue!7} HyperGraphRAG~\citep{HyperGraphRAG} & 36.0 & - & 34.2 & - & - & - & 34.6 & 34.2 & 35.7 & 34.3 & 34.5 & 35.8 & 35.7 & 35.4 & 34.8 & 36.1 & 11.7 & \textcolor{red}{-32.7} \\
\rowcolor{teal!8} iText2KG~\citep{iText2KG} & 33.5 & - & 33.0 & - & 40.6 & - & 34.0 & 35.1 & 34.1 & 34.6 & 34.1 & 36.0 & 37.5 & 37.8 & 36.8 & 37.2 & 17.8 & \textcolor{red}{-26.6} \\
\rowcolor{teal!8} OntoKG~\citep{OntoKG} & 34.4 & - & 33.2 & - & - & - & 33.2 & 33.3 & 34.2 & 32.8 & 33.2 & 34.5 & 34.3 & 34.4 & 33.5 & 34.6 & 11.3 & \textcolor{red}{-33.1} \\
\rowcolor{teal!14} KAG/OpenSPG~\citep{KAG} & 36.5 & - & 34.4 & - & 42.8 & - & 36.6 & 36.8 & 36.5 & 37.0 & 36.9 & 38.0 & 39.5 & 39.7 & 38.7 & 39.5 & 18.9 & \textcolor{red}{-25.5} \\
\rowcolor{teal!14} Graphiti~\citep{Graphiti} & 34.5 & - & 34.0 & - & 44.2 & 54.4 & 38.9 & 41.1 & 41.0 & 44.3 & 43.5 & 42.1 & 39.9 & 46.3 & 39.6 & 40.9 & 27.8 & \textcolor{red}{-16.6} \\
\rowcolor{orange!8} MemGPT/Letta~\citep{MemGPT} & 25.2 & - & 31.8 & - & 41.8 & - & 31.1 & 33.0 & 31.5 & 32.4 & 31.6 & 33.2 & 34.3 & 34.4 & 33.9 & 34.0 & 16.5 & \textcolor{red}{-27.9} \\
\rowcolor{orange!8} Mem0~\citep{Mem0} & 23.6 & - & 31.0 & - & 41.1 & - & 30.6 & 31.6 & 29.8 & 31.2 & 30.8 & 31.7 & 33.3 & 34.0 & 32.8 & 33.2 & 16.0 & \textcolor{red}{-28.4} \\
\rowcolor{orange!8} A-MEM~\citep{A-MEM} & 28.3 & - & 33.0 & - & 42.9 & 50.7 & 35.7 & 37.2 & 38.0 & 41.0 & 40.4 & 39.0 & 37.7 & 43.7 & 36.6 & 37.6 & 25.8 & \textcolor{red}{-18.6} \\
\rowcolor{yellow!12} OmniRAG-Agent~\citep{OmniRAG-Agent} & 51.7 & 63.3 & 32.4 & 46.7 & - & - & 51.1 & 46.7 & 47.6 & 49.5 & 45.0 & 49.1 & 50.3 & 48.4 & 48.3 & 49.2 & 32.4 & \textcolor{red}{-12.0} \\
\rowcolor{green!8} \textbf{OpenClaw}~\citep{OpenClaw} & 51.0 & 62.2 & 27.3 & 43.0 & 38.5 & 44.5 & 42.8 & 41.3 & 42.6 & 47.1 & 43.8 & 45.5 & 44.6 & 47.9 & 43.6 & 44.9 & 44.4 & \textbf{0.0} \\
\rowcolor{green!8} Cognee~\citep{Cognee} & 52.6 & 64.0 & 28.5 & 43.9 & 38.9 & 45.1 & 44.4 & 43.7 & 43.5 & 47.7 & 44.2 & 46.4 & 45.7 & 49.3 & 43.9 & 46.3 & 45.5 & \textcolor{green!60!black}{+1.1} \\
\rowcolor{green!8} LLM Wiki~\citep{LLMWiki} & 52.9 & 64.5 & 28.9 & 44.2 & 38.9 & 47.7 & 45.1 & 43.9 & 44.2 & 48.0 & 45.4 & 47.1 & 46.8 & 49.9 & 45.1 & 46.2 & 46.2 & \textcolor{green!60!black}{+1.8} \\
\rowcolor{green!8} Graphify~\citep{Graphify} & 52.2 & 63.5 & 28.5 & 43.7 & 39.0 & 44.9 & 44.7 & 42.5 & 43.5 & 48.1 & 44.4 & 46.3 & 45.0 & 48.9 & 43.9 & 45.7 & 45.3 & \textcolor{green!60!black}{+0.9} \\
\midrule
\rowcolor{red!8} KC w/o C1 & 39.0 & - & 34.4 & - & 45.1 & 58.0 & 41.9 & 43.3 & 44.7 & 46.2 & 45.4 & 44.4 & 43.5 & 47.7 & 41.8 & 42.7 & 29.4 & \textcolor{red}{-15.0} \\
\rowcolor{red!8} KC w/o C2 & 74.9 & 89.3 & 34.5 & 48.9 & - & 64.3 & 58.0 & 54.9 & 57.4 & 59.2 & 55.1 & 56.1 & 57.3 & 57.3 & 54.9 & 56.5 & 52.0 & \textcolor{green!60!black}{+7.6} \\
\rowcolor{red!8} KC w/o C3 & 74.9 & 89.3 & 24.0 & 35.5 & 39.9 & 48.1 & 52.7 & 50.7 & 52.0 & 53.6 & 49.8 & 51.6 & 52.9 & 53.3 & 50.9 & 52.0 & 51.9 & \textcolor{green!60!black}{+7.5} \\
\rowcolor{red!12} \textbf{KnowledgeClaw} & \textbf{74.9} & \textbf{89.3} & \textbf{34.5} & \textbf{48.9} & \textbf{46.8} & \textbf{64.3} & \textbf{60.6} & \textbf{58.2} & \textbf{60.2} & \textbf{62.0} & \textbf{57.9} & \textbf{59.1} & \textbf{60.8} & \textbf{61.0} & \textbf{58.2} & \textbf{59.7} & \textbf{59.8} & \textcolor{green!60!black}{\textbf{+15.4}} \\
\bottomrule
\end{tabular}%
}
\end{table}

\subsection{Main Results (RQ1)}

\textbf{Overall.} KnowledgeClaw is the only system to score on all six OmniKE tasks and ranks first on every one, reaching aggregate OKE 59.8 and beating the strongest external baseline LLM Wiki by +13.6 and our prior anchor OpenClaw by +15.4; baselines covering only one or two of properties (i)--(iii) cluster below 50, text-only and memory-only systems stay under 30, and the four general agentic KE systems land in 44--47, escaping the text-only trap yet still missing at least one axis.

\textbf{Per-task.} On the \emph{structural} tasks T1/T2, KC reaches 74.9/89.3 (+22.0/+24.6 over the best baseline), confirming C1's hypergraph captures richer cross-modal structure than text-only extraction; on the \emph{dynamic} tasks T5/T6, KC attains 46.8/64.3 (+2.6/+6.1), validating both C2's lifecycle loop and C3's community synthesis; on the \emph{retrieval} tasks T3/T4, KC attains 34.5/48.9 with smaller but consistent gains, since the multi-granularity payoff is bounded by the shared backbone.

\textbf{Per-persona.} Gains are stable across all ten personas, with KC's per-persona OKE in 58.2--62.0---the smallest variance in the table; the largest advantages appear in modality-heavy domains (Healthcare, Law, Robotics, Building, all $>$+13 OKE), and even text-centric Software, Education, and Finance retain stable +10--12 gains.

\subsection{Ablation Study (RQ2)}

Removing each contribution in turn (last four rows of Table~\ref{tab:main}) hurts precisely the axis it serves. \textbf{w/o C1} collapses cross-modal coreference: T2 and T4 become unscorable and OKE falls to 29.4 ($-$30.4)---by far the largest single-component drop and clear evidence that C1 is the structural foundation. \textbf{w/o C2} loses lifecycle: T5 becomes unscorable and OKE drops to 52.0 ($-$7.8) with the other five tasks preserved, showing C2 targets lifecycle alone. \textbf{w/o C3} collapses retrieval to a flat top-$k$ pass: T3 drops 34.5$\to$24.0, T6 drops 64.3$\to$48.1, OKE falls to 51.9 ($-$7.9), while structural T1/T2 remain unchanged. This orthogonal decay pattern validates our central claim: omnimodality, closed-loop lifecycle, and multi-angle application form a minimum-sufficient set of axes mapped one-to-one onto C1--C3.

\subsection{Analysis of Modality and Persona Generalization (RQ3)}\label{5.4}

Figure~\ref{fig:gen} decomposes KnowledgeClaw's gain by persona and modality, testing whether the omnimodal hypergraph yields uniform wins across domains and whether each modality is needed.

\begin{figure}[h]
\centering
\includegraphics[width=\textwidth]{../plot/RQ3/fig_rq3.pdf}
\caption{(a) Persona$\times$Task heatmap of $\mathrm{KC}-\mathrm{max(baseline)}$ (positive = KC wins). (b) Score drop on each task when one modality is removed from $\mathcal{G}$ (larger = more load-bearing).}\label{fig:gen}
\end{figure}

\textbf{Every modality earns its place.} Panel (b) reads as ``what breaks if a modality disappears.'' Removing audio or video causes the largest drops on T2 (cross-modal coreference) and T4 (cross-modal reasoning), while image and PDF dominate the structural signal on T1 and T3; no single modality can be substituted by text, so the omnimodal substrate is load-bearing rather than decorative.

\textbf{Gains are uniform across personas.} Panel (a) is positive on virtually every persona$\times$task cell: KC beats the strongest baseline on each of the ten personas, with the biggest deltas in modality-heavy domains (Healthcare, Robotics, Building) and a stable +10--12 cushion even on text-centric Software and Finance, indicating the gain comes from the design rather than a favorable domain mix.

\subsection{Analysis of Closed-Loop Lifecycle Robustness (RQ4)}\label{5.5}

Figure~\ref{fig:cycle} stress-tests C2 by streaming 200 Add/Delete/Refresh events from each persona's audit log, asking whether the live hypergraph stays consistent or decays as edits accumulate.

\begin{figure}[h]
\centering
\includegraphics[width=\textwidth]{../plot/RQ4/fig_rq4.pdf}
\caption{(a) T5 consistency over a 200-event CRUD stream for six methods (mean $\pm$ 1$\sigma$ over 5 seeds). (b) Per-method, per-event-type average score impact for Add/Delete/Refresh.}\label{fig:cycle}
\end{figure}

\textbf{No drift over a 200-event stream.} KC's T5(t) curve in panel (a) stays essentially flat across the full horizon, while the strongest CRUD-capable baselines (Graphiti, A-MEM) drift downward by over 10 points and KC w/o C2 collapses outright. The cascading delete and in-place refresh in Eqs.~\eqref{eq:cascade}--\eqref{eq:refresh} clear orphans and dangling hyperedges before they can poison retrieval, so answer fidelity tracks the corpus instead of slowly decaying.

\textbf{Robust on every edit type.} Panel (b) shows baselines lose the most on \textsc{Delete} (through orphan accumulation) and \textsc{Refresh} (through stale facts), while KC's per-event-type impact stays uniformly small. The Build--Query--Reason--Update loop therefore handles all three operators with comparable robustness, rather than over-fitting to any single edit type.

\clearpage

\subsection{Analysis of Intent Routing Effectiveness (RQ5)}\label{5.6}

Figure~\ref{fig:route} compares the lightweight LLM router against an oracle and three single-angle policies on four query types, asking whether routing is needed and whether one LLM call suffices.

\begin{figure}[h]
\centering
\includegraphics[width=\textwidth]{../plot/RQ5/fig_rq5.pdf}
\caption{Left: 3$\times$3 router confusion matrix (oracle row vs.\ predicted column; cell = recall). Right: per-query-type score for sub-only, com-only, disc-only, oracle, and the KC router.}\label{fig:route}
\end{figure}

\textbf{One LLM call is enough.} The confusion matrix is strongly diagonal: the router recovers the oracle angle on most queries (only \textsc{com}$\leftrightarrow$\textsc{disc} is mildly confused), so a single LLM call replaces a multi-pass classifier.

\textbf{No single angle suffices.} sub-only loses on synthesis/discovery, com-only on multi-hop, disc-only outside discovery; the KC router stays within 1--2 points of the oracle on every type, dominating any fixed retrieval mode.

\subsection{Cost and Time Analysis (RQ6)}\label{5.7}

Figure~\ref{fig:cost} reports the two costs a deployment user actually pays: dollars spent per persona on the OmniKE pipeline (left) and wall-clock construction time on the full-coverage tier (right).

\begin{figure}[h]
\centering
\includegraphics[width=\textwidth]{../plot/RQ6/fig_rq6.pdf}
\caption{Left: OKE vs.\ USD spend per persona (log scale); KC (starred) sits on the Pareto frontier. Right: construction time per persona on T1--T6 for the five full-coverage methods.}\label{fig:cost}
\end{figure}

\textbf{Cheaper.} KC spends \$1.95 per persona, \textbf{20--34\% below} OpenClaw (\$2.45), Cognee (\$2.60), LLM Wiki (\$2.75), and Graphify (\$2.95): C1 reuses one LLM pass per asset, C2 updates incrementally, and C3 routes only to the relevant angle.

\textbf{Faster.} On the same tier KC finishes T1--T6 in 41\,min---\textbf{14\,min} ahead of OpenClaw and \textbf{35\,min} ahead of Graphify---with 580\,ms serving latency, making KC cheaper, faster, and +13.6 OKE higher than every full-coverage baseline.

\section{Conclusion}\label{sec:conclusion}

\textbf{KnowledgeClaw} is the first agentic knowledge engineering system that jointly satisfies omnimodality, closed-loop lifecycle, and multi-angle application on a unified omnimodal hypergraph; on OmniKE-Bench against 22 baselines it ranks first on all six tasks (OKE 59.8, +13.6 over the strongest) at lower cost and shorter time, catalyzing further research on live all-modality knowledge engineering.

\clearpage

\section{Related Work}

\subsection{Knowledge Engineering}

Knowledge engineering began with expert systems and ontologies~\citep{KEclassic, CYC, WordNet}, which required manual encoding of concepts and inference rules. Web-scale corpora then shifted the paradigm to automated knowledge graph (KG) construction~\citep{Freebase, DBpedia, YAGO, Wikidata}, maturing along extraction~\citep{NELL, OpenIE}, embedding~\citep{TransE, ComplexE}, and linking~\citep{BLINK, KG}. LLMs have since rewritten this paradigm~\citep{KGLLMsurvey}, both by directly extracting and normalizing entities and relations~\citep{LLMIEsurvey, KAG, OntoKG, iText2KG} and by injecting external knowledge into generation via RAG~\citep{RAG, RAGSurvey}. Centered on graph structure, this trajectory spans hierarchical summarization~\citep{RAPTOR}, community-level synthesis~\citep{GraphRAG, LightRAG}, path pruning and multi-hop/proposition reasoning~\citep{PathRAG, HopRAG, PropRAG, E2GraphRAG}, heterogeneous and linearized graphs~\citep{NodeRAG, LinearRAG}, causal and domain-specialized graphs~\citep{CausalRAG, MedGraphRAG, PIKE-RAG}, hypergraph modeling of higher-order relations~\citep{HyperGraphRAG, HyperRAG, EHRAG}, and graph foundation models~\citep{GFM-RAG, G-reasoner}. On top of this, non-parametric continual learning~\citep{HippoRAG2}, KG-based reasoning scaffolds~\citep{ToG}, and agentic retrieval~\citep{KBQA-o1, GraphSearch, Youtu-GraphRAG, AgenticRAGSurvey, Graph-R1} push GraphRAG toward a learnable and interactive regime. Yet these methods remain confined to single-modality text or static-corpus QA, without joint treatment of omnimodality, closed-loop lifecycle, and multi-facet application.

\subsection{LLMs and Agents}

Built on the Transformer~\citep{Transformer} and the pre-training--RLHF paradigm~\citep{BERT, GPT3, InstructGPT}, modern LLMs have scaled with Chain-of-Thought reasoning~\citep{CoT, ToT, LLMsurvey} and, through large-scale reinforcement learning, evolved into the ``reasoning model'' paradigm~\citep{GPT-o1, DeepSeek-R1, Qwen3}. On the modality axis, progress has moved from vision--language and speech alignment~\citep{CLIP, Whisper, Llava} to production-grade VLMs~\citep{GPT4, Qwen3-VL}, and further to omnimodal backbones unifying text, image, audio, and video~\citep{Qwen3-Omni, Qwen3.5-Omni, Dynin-Omni, OmniRAG-Agent, OmniGAIA, Gemini3}. LLM-based agents have matured rapidly~\citep{AgentSurvey}: the reason--act--retrieve--reflect paradigm~\citep{ReAct, Reflexion, SelfRAG, Toolformer} laid the capability foundation, and a tighter coupling of reasoning, tool use, and retrieval yielded ``deep-research''-style agents~\citep{Search-o1, Search-R1}. Agents have also enteredW real digital environments for web and desktop operation~\citep{WebArena, OSWorld, ClawGUI}, acquired CRUD interfaces and long/short-term memory via memory-oriented architectures~\citep{Graphiti, MemGPT, Mem0, A-MEM, AgenticMemory, HyperMem}, and achieved end-to-end autonomy on structured tasks such as software engineering~\citep{SWE-agent, OpenHands, ClaudeCode}. Recently, OpenClaw~\citep{OpenClaw} unifies perception, reasoning, tool use, and memory in a single closed-loop agent, opening a promising path toward the omnimodal knowledge engineering task~\citep{Cognee, LLMWiki, Graphify} targeted in this work.

\clearpage



\bibliography{custom}{}
\bibliographystyle{plain}


%%%%%%%%%%%%%%%%%%%%%%%%%%%%%%%%%%%%%%%%%%%%%%%%%%%%%%%%%%%%

\newpage
\appendix
\section*{Appendix} 

\section{KnowledgeClaw Prompts}\label{app:prompts}

This appendix lists every LLM prompt invoked by KnowledgeClaw, mirroring the prompt-card presentation of HyperGraphRAG~\citep{HyperGraphRAG}. Prompts are grouped by the contribution they serve --- \textbf{C1 Construction} (\S\ref{sec:c1}), \textbf{C2 Closed-loop Lifecycle} (\S\ref{sec:c2}), and \textbf{C3 Multi-angle Application} (\S\ref{sec:c3}) --- and are reproduced verbatim from the released code, with \texttt{\{x\}} denoting a runtime placeholder. The user-message wrapper around system prompts is kept compact since the system prompt already constrains the format.

\subsection{Construction Prompts (C1)}

\noindent{\textbf{Prompt P1 --- Entity and Proposition Extraction} ($\mathrm{Extract}_{\mathrm{LLM}}$, T1).} For each chunk produced from a text/PDF/image-caption/transcript asset, the extractor emits atomic propositions with their grounded entity mentions; canonical UPPER-CASED names feed the within-modality canonicalization $\kappa$ and stabilize lexical variants across documents.
\begin{shaded}
\footnotesize
\begin{verbatim}
Extract key entities and propositions from the text below.
For each sentence/proposition identify the named entities (people,
organisations, concepts, methods, tools, places).
For each entity provide:
  - "name": the canonical English name, UPPER-CASED. Resolve abbreviations
    and surface variants to a single canonical form.

Return JSON ONLY (no markdown, no explanation):
{
  "propositions": [
    { "sentence": "atomic fact or claim",
      "entities": [ {"name": "ENTITY_NAME"} ] }
  ]
}
Text: {text}
\end{verbatim}
\end{shaded}

\subsection{Closed-loop Lifecycle Prompts (C2)}

\noindent{\textbf{Prompt P2 --- Lifecycle Health Audit} (used during the \textsc{Update} stage of Eq.~\ref{eq:loop}).} After every CRUD pass, the lifecycle agent runs a deterministic structural summary over $\mathcal{G}$ (entity/hyperedge/asset/cross-modal counts, isolation ratio, top entities by degree) and asks the LLM for a short qualitative read on coverage, connectivity, and cross-modal integration. The output is logged to $\mathcal{L}_n$ and used to flag silent decay before it reaches retrieval.
\begin{shaded}
\footnotesize
\begin{verbatim}
Knowledge Base Statistics:
- Total nodes: {total_nodes}
- Entity nodes / Hyperedge nodes / Cross-modal hyperedges: {ent}/{he}/{xm}
- Asset nodes / Total edges / Chunks: {asset}/{edges}/{chunks}
- Isolated nodes: {iso}        - Average degree: {avg_deg}
- Top entities (by degree): {top_entities}

Write a 2-3 sentence health assessment of this typed omnimodal knowledge
hypergraph. Comment on coverage, connectivity, cross-modal integration,
and any concerns.
\end{verbatim}
\end{shaded}

\subsection{Multi-angle Application Prompts (C3)}

\noindent{\textbf{Prompt P3 --- Community Report Synthesis} ($R_{\mathrm{com}}$, Eq.~\ref{eq:com}).} After Louvain partitions the cooccurrence graph, each community $\mathcal{C}_k$ is summarized into $\sigma_k$, embedded, and stored in a dedicated VDB so that topic queries can dispatch to the right community in a single cosine probe instead of re-traversing $\mathcal{G}$.
\begin{shaded}
\footnotesize
\begin{verbatim}
You are summarising one community of entities from a personal knowledge
hypergraph.

Community entities (canonical IDs): {entities}

Sample propositions / hyperedges that mention these entities:
{propositions}

Modalities present in this community: {modalities}

Write a concise community report with the following structure (no markdown
headers, just plain prose):
1. THEME — a single-sentence label for what this community is about.
2. KEY ENTITIES — list the 3-5 most central entities and their role.
3. SUMMARY — a 3-5 sentence synthesis of the relationships and facts.
4. CROSS-MODAL NOTES — if multiple modalities appear, note how they
   reinforce each other.

Keep the whole report under 200 words and avoid generic boilerplate.
\end{verbatim}
\end{shaded}

\noindent{\textbf{Prompt P4 --- Intent Router} ($\pi^{*}$, Eq.~\ref{eq:disc-route}).} A keyword heuristic resolves the obvious factual/synthesis extremes for free; only ambiguous queries fall through to the LLM router below, which classifies into the three angles in $\mathcal{R}$ at $\leq\!8$ output tokens.
\begin{shaded}
\footnotesize
\begin{verbatim}
Classify the following user query into ONE label that describes the BEST
retrieval strategy over a personal knowledge hypergraph.

Labels:
  - subgraph  : precise factual lookup, single fact, definition,
                who/when/where/which
  - community : open-ended synthesis, summarisation,
                "what do I know about X", thematic overviews
  - hybrid    : ambiguous or multi-hop questions that need both precise
                facts and broader context

Reply with EXACTLY one of: subgraph, community, hybrid

Query: {query}
Label:
\end{verbatim}
\end{shaded}

\noindent{\textbf{Prompt P5 --- Subgraph + Community QA} ($\mathrm{LLM}_{\mathrm{gen}}$, Eq.~\ref{eq:disc-route}; T3, T4).} The QA agent receives the routed evidence as a concatenation $\bigoplus_{a\in\pi^{*}(q)} R_a(q)$ where each chunk/community block carries a citation tag, allowing the system to compute T4's cross-modal evidence rate from the cited modalities.
\begin{shaded}
\footnotesize
\begin{verbatim}
[system]
You are a precise knowledge assistant for a personal knowledge hypergraph.
Answer the user's question using ONLY the provided evidence. If the evidence
is insufficient, say so clearly. Cite chunk-N or community-N when you
reference a fact.

[user]
Question: {question}
Routed intent: {intent}
Evidence:
{evidence}
Answer:
\end{verbatim}
\end{shaded}

\noindent{\textbf{Prompt P6 --- ReAct Inference Agent} (multi-round $R_{\mathrm{sub}}\!\cup\!R_{\mathrm{com}}$).} For deep-reasoning queries that exceed a single retrieval pass: each \texttt{<search>} tag re-enters the hierarchical retriever via the same router as P4, and a final \texttt{<insight>} tag closes the trace after at most \texttt{\{max\_rounds\}} rounds.
\begin{shaded}
\footnotesize
\begin{verbatim}
You are a knowledge inference engine analysing a personal knowledge
hypergraph.

You have a search tool. Use it to find relevant knowledge, then synthesise
deep insights.

Format your response with these XML-like tags:
<think>your reasoning about what to search next or what you have learned
so far</think>
<search>a specific search query to retrieve relevant knowledge</search>

After you have gathered enough knowledge (usually 3-5 searches), produce
your final insight:
<insight>your synthesised insight, prediction, or conclusion — cite
specific facts from the knowledge base</insight>

Rules:
- Each <search> triggers an actual KB retrieval through the hierarchical
  retriever (subgraph + community).
- You may do up to {max_rounds} search rounds before you must produce
  <insight>.
- Always base insights on retrieved facts; note when evidence is sparse.
- Write the <insight> tag LAST, after all your thinking and searching.
\end{verbatim}
\end{shaded}

\noindent{\textbf{Prompt P7 --- Knowledge Discovery / Gap Analysis} ($\mathrm{LLM}_{\mathrm{disc}}$, T6).} Realizes the discovery triplet $(\mathrm{known}(t),\mathrm{gaps}(t),\mathrm{sug}(t))$ from Eq.~\ref{eq:disc-route}: $\mathrm{known}$ is the cosine threshold $\theta_{\mathrm{k}}$ over $\phi(\cdot)$, while $\mathrm{gaps},\mathrm{sug}$ are produced by the LLM under topic-conditioned community context $\{\sigma_k\}_{\mathcal{C}_t^{\mathrm{com}}}$.
\begin{shaded}
\footnotesize
\begin{verbatim}
You are a research analyst for a personal knowledge hypergraph.

Given the topic and the retrieved knowledge snippets (which include both
fine-grained chunks and community-level summaries), you must:
1. Summarise what the KB already knows about the topic.
2. Identify specific knowledge gaps — what is missing or unclear.
3. Suggest 5-10 concrete search queries or sub-topics to fill those gaps.
4. Rate the current coverage: low / medium / high.

Return a structured analysis.
\end{verbatim}
\end{shaded}

\textbf{Reproducibility.} The seven prompts above, together with the deterministic operators in Eqs.~\eqref{eq:hg}--\eqref{eq:disc-route}, fully specify KnowledgeClaw's LLM surface; all run at \texttt{temperature}\,$\in\![0,0.4]$ on the unified \texttt{gemini-3.1-flash-lite-preview} backbone (\S\ref{sec:exp}).

\clearpage

\section{Proof of Propositions}\label{proofs}

This appendix proves Propositions~1--3 stated in \S\ref{sec:method}, formalising the three claims that ground C1 (omnimodal lossless representation), C2 (lifecycle consistency under finite event streams), and C3 (necessity of multi-angle routing). Notation follows \S\ref{sec:c1}--\S\ref{sec:c3}; we write $H(\cdot)$ for Shannon entropy and $I(\cdot;\cdot)$ for mutual information.

\subsection{Proof of Proposition 1}\label{proof1}
\textbf{Proposition 1.} \textit{$\mathcal{G}$ strictly subsumes text-only graphs: $n$-ary cross-modal facts on any $\mathcal{M}'\!\subseteq\!\mathcal{M}$ are uniquely representable in $\mathcal{G}$ yet unrecoverable from text, generalizing uniformly across modalities.}\vspace{-2mm}
\begin{proof}
Let a fact be denoted as
\begin{equation}
F\;=\;(\mathcal{E}_F,\,\mathcal{A}_F,\,\pi_F,\,\tau_F,\,\omega_F),\qquad |\mathcal{E}_F|\!=\!n\!\geq\!2,\qquad \mathcal{M}_F\!=\!\mu(\mathcal{A}_F)\!\subseteq\!\mathcal{M}.
\label{eq:p1-fact}
\end{equation}
The omnimodal encoder of Eq.~\eqref{eq:hg} sends $F$ to a typed hyperedge,
\begin{equation}
\varphi_{\mathcal{G}}\colon F\;\longmapsto\;h\!=\!(\mathcal{E}_F,\,\tau_F,\,\pi_F,\,\mathcal{A}_F,\,\omega_F)\;\in\;\mathcal{H},
\label{eq:p1-phiG}
\end{equation}
while any text-only graph $\mathcal{G}_T\!=\!(\mathcal{E}_T,E_T)$ admits at most the projection
\begin{equation}
\varphi_T\colon F\;\longmapsto\;\bigl(\mathcal{E}_F\!\cap\!\mathcal{E}_T,\;\{(u,v)\!:\!u,v\!\in\!\mathcal{E}_F\!\cap\!\mathcal{E}_T,\,u\!\neq\!v,\,\exists f\!\in\!\mathcal{A}_F\,m(f)\!=\!\textit{text}\}\bigr).
\label{eq:p1-phiT}
\end{equation}

Since the tuple $(\mathcal{E}_h,\mathcal{A}_h,\pi_h)$ uniquely recovers $\mathcal{M}_h\!=\!\mu(\mathcal{A}_h)$ and $F$, the map $\varphi_{\mathcal{G}}$ is injective, so
\begin{equation}
H\!\bigl(F\!\mid\!\varphi_{\mathcal{G}}(F)\bigr)\;=\;0,\qquad I\!\bigl(F;\varphi_{\mathcal{G}}(F)\bigr)\;=\;H(F).
\label{eq:p1-Igamma}
\end{equation}

For $\varphi_T$, fix any cross-modal subset $\mathcal{M}'\!\subseteq\!\mathcal{M}$ with $|\mathcal{M}'|\!\geq\!2$ and pick $F_1$ with $\mathcal{M}_{F_1}\!=\!\mathcal{M}'$. Pick a non-text witness $f^{*}\!\in\!\mathcal{A}_{F_1}\!\setminus\!\mathcal{A}_{F_1}^{\textit{text}}$ with $m(f^{*})\!\in\!\mathcal{M}'\!\setminus\!\{\textit{text}\}$, and construct
\begin{equation}
F_2\;=\;\bigl(\mathcal{E}_{F_1},\;(\mathcal{A}_{F_1}\!\setminus\!\{f^{*}\})\!\cup\!\{\bar{f}^{*}\},\;\bar{\pi}_{F_1},\;\tau_{F_1},\;1\!-\!\omega_{F_1}\bigr),\qquad m(\bar{f}^{*})\!=\!m(f^{*}),\;\omega(\bar{f}^{*})\!=\!1\!-\!\omega(f^{*}).
\label{eq:p1-construct}
\end{equation}
Since the edited asset and proposition lie outside the text projection, both facts collapse to the same image:
\begin{equation}
\varphi_T(F_1)\;=\;\bigl(\mathcal{E}_{F_1}\!\cap\!\mathcal{E}_T,\,E_T(\mathcal{A}_{F_1}^{\textit{text}})\bigr)\;=\;\bigl(\mathcal{E}_{F_2}\!\cap\!\mathcal{E}_T,\,E_T(\mathcal{A}_{F_2}^{\textit{text}})\bigr)\;=\;\varphi_T(F_2)\;=:\;g,
\label{eq:p1-collide-img}
\end{equation}
while $F_1\!\neq\!F_2$ in $\mathcal{E}\!\times\!2^{\mathcal{D}}$ because $\bar{f}^{*}\!\neq\!f^{*}$, so the preimage carries strictly more than one element:
\begin{equation}
\{F_1,F_2\}\;\subseteq\;\varphi_T^{-1}(g)\quad\Longrightarrow\quad\bigl|\varphi_T^{-1}(g)\bigr|\!\geq\!2,\qquad p(F_1\!\mid\!\varphi_T(F)\!=\!g),\,p(F_2\!\mid\!\varphi_T(F)\!=\!g)\!>\!0.
\label{eq:p1-collide}
\end{equation}

The conditional distribution $p(F\!\mid\!\varphi_T(F)\!=\!g)$ therefore has at least two atoms with positive mass, giving a strictly positive point-wise conditional entropy:
\begin{equation}
H\!\bigl(F\!\mid\!\varphi_T(F)\!=\!g\bigr)\;=\;-\!\!\sum_{F_i\in\varphi_T^{-1}(g)}\!\!\frac{p(F_i)}{p(g)}\,\log_2\!\frac{p(F_i)}{p(g)}\;>\;0.
\label{eq:p1-Hpoint}
\end{equation}
Marginalising over the projection $y$ bounds the unconditional conditional entropy from below,
\begin{equation}
H\!\bigl(F\!\mid\!\varphi_T(F)\bigr)\;=\;\sum_{y}p(y)\,H\!\bigl(F\!\mid\!\varphi_T(F)\!=\!y\bigr)\;\geq\;p(g)\,H\!\bigl(F\!\mid\!\varphi_T(F)\!=\!g\bigr)\;>\;0,
\label{eq:p1-Hmarg}
\end{equation}
which by definition translates into a strict mutual-information loss:
\begin{equation}
I\!\bigl(F;\varphi_T(F)\bigr)\;=\;H(F)\;-\;H\!\bigl(F\!\mid\!\varphi_T(F)\bigr)\;<\;H(F).
\label{eq:p1-Iloss}
\end{equation}

Combining Eq.~\eqref{eq:p1-Igamma} with Eq.~\eqref{eq:p1-Iloss} yields the strict information gap
\begin{equation}
I\!\bigl(F;\varphi_{\mathcal{G}}(F)\bigr)\;-\;I\!\bigl(F;\varphi_T(F)\bigr)\;=\;H\!\bigl(F\!\mid\!\varphi_T(F)\bigr)\;>\;0,\qquad \forall\,\mathcal{M}'\!\subseteq\!\mathcal{M},\;|\mathcal{M}'|\!\geq\!2.
\label{eq:p1-gap}
\end{equation}
The construction depends only on $|\mathcal{M}_F|\!\geq\!2$, so Eq.~\eqref{eq:p1-gap} holds uniformly across the modality lattice; in the degenerate single-modality case $\mathcal{M}_F\!=\!\{\textit{text}\},\,n\!=\!2$ the two encoders coincide and Eq.~\eqref{eq:p1-gap} reduces to equality. Hence $\mathcal{G}$ strictly subsumes any text-only graph on every cross-modal subset of $\mathcal{M}$, and Proposition~1 is established.
\end{proof}

\subsection{Proof of Proposition 2}\label{proof2}
\textbf{Proposition 2.} \textit{Under any finite event sequence, the lifecycle agent preserves consistency: $\mathcal{S}(e)\!\neq\!\emptyset$ for every $e\!\in\!\mathcal{E}_N$ and $\mathcal{E}_h\!\subseteq\!\mathcal{E}_N$ for every $h\!\in\!\mathcal{H}_N$.}\vspace{-2mm}
\begin{proof}
Define two structural invariants on the lifecycle state $\mathcal{G}_n$:
\begin{equation}
\text{(I1)}\quad\forall\,e\!\in\!\mathcal{E}_n\!:\;\mathcal{S}(e)\!\neq\!\emptyset,\qquad
\text{(I2)}\quad\forall\,h\!\in\!\mathcal{H}_n\!:\;\mathcal{E}_h\!\subseteq\!\mathcal{E}_n.
\label{eq:p2-inv}
\end{equation}
We argue by induction on $n\!=\!0,\dots,N$, where $\mathcal{G}_n$ evolves under the lifecycle operator
\begin{equation}
\mathcal{G}_{n+1}\;=\;(\mathrm{Update}\!\circ\!\mathrm{Reason}\!\circ\!\mathrm{Query}\!\circ\!\mathrm{Build})(\mathcal{G}_n,\Delta_n),\qquad \Delta_n\!\in\!\{\textsc{Add},\textsc{Del},\textsc{Refresh}\}\!\times\!\mathcal{D},
\label{eq:p2-loop}
\end{equation}
of Eq.~\eqref{eq:loop}. The base case is immediate: $\mathcal{G}_0\!=\!(\emptyset,\emptyset,\emptyset,\emptyset)$ trivially satisfies (I1)--(I2).

For the inductive step, assume (I1)--(I2) at $\mathcal{G}_n$. When $\Delta_n\!=\!\textsc{Add}(f)$, the four sub-operators of Eq.~\eqref{eq:p2-loop} act as
\begin{equation}
\begin{aligned}
&\text{Build}:\;(\hat{\mathcal{E}}_f,\hat{\mathcal{P}}_f,\hat{\mathcal{H}}_f,\hat{\mu}_f)\;=\;\mathrm{Extract}_{\mathrm{LLM}}\!\bigl(f,\,m(f)\bigr),\;\;\forall e\!\in\!\hat{\mathcal{E}}_f\!:\;f\!\in\!\mathcal{S}(e),\;\;\forall h\!\in\!\hat{\mathcal{H}}_f\!:\;\mathcal{E}_h\!\subseteq\!\hat{\mathcal{E}}_f\!\cup\!\mathcal{E}_n,\\[2pt]
&\text{Query/Reason}:\;\hat{\chi}_f\;=\;\bigl\{(e_a,e_b)\!:\!e_a\!\in\!\hat{\mathcal{E}}_f,\,e_b\!\in\!\mathcal{E}_n,\,\chi(e_a,e_b)\!=\!1\bigr\}\;\subseteq\;\mathcal{H},\\[2pt]
&\text{Update}:\;\mathcal{E}_{n+1}\;=\;\mathcal{E}_n\!\cup\!\kappa(\hat{\mathcal{E}}_f),\qquad \mathcal{H}_{n+1}\;=\;\mathcal{H}_n\!\cup\!\hat{\mathcal{H}}_f\!\cup\!\hat{\chi}_f,\qquad \mathcal{A}_{n+1}\;=\;\mathcal{A}_n\!\cup\!\{f\}.
\end{aligned}
\label{eq:p2-add-state}
\end{equation}
The canonicalisation $\kappa$ only merges within-modality variants, so $\mathcal{S}(\kappa(e))\!\supseteq\!\mathcal{S}(e)$ for every merged $e$, and the coreference edges in $\hat{\chi}_f$ are binary hyperedges over already-existing entities. Both (I1) and (I2) therefore carry over: every $e\!\in\!\mathcal{E}_{n+1}$ either inherits $\mathcal{S}(e)\!\neq\!\emptyset$ from $\mathcal{G}_n$ or has $f\!\in\!\mathcal{S}(e)$, and every $h\!\in\!\mathcal{H}_{n+1}$ has $\mathcal{E}_h\!\subseteq\!\mathcal{E}_n\!\cup\!\hat{\mathcal{E}}_f\!\subseteq\!\mathcal{E}_{n+1}$.

When $\Delta_n\!=\!\textsc{Del}(f)$, the cascading-delete operator of Eq.~\eqref{eq:cascade} first identifies the orphan set
\begin{equation}
\mathcal{E}_n^{\mathrm{orph}}(f)\;=\;\bigl\{\,e\!\in\!\mathcal{E}_n\,:\,\mathcal{S}(e)\!\subseteq\!\{f\}\,\bigr\},
\label{eq:p2-orph}
\end{equation}
then evicts orphans, supporting asset, and dangling hyperedges in one transactional pass:
\begin{equation}
\mathcal{E}_{n+1}\!=\!\mathcal{E}_n\!\setminus\!\mathcal{E}_n^{\mathrm{orph}}(f),\;\;
\mathcal{H}_{n+1}\!=\!\bigl\{\,h\!\in\!\mathcal{H}_n\!:\,\mathcal{A}_h\!\not\subseteq\!\{f\}\;\wedge\;\mathcal{E}_h\!\cap\!\mathcal{E}_n^{\mathrm{orph}}(f)\!=\!\emptyset\,\bigr\},\;\;
\mathcal{A}_{n+1}\!=\!\mathcal{A}_n\!\setminus\!\{f\}.
\label{eq:p2-del-state}
\end{equation}
For every $e\!\in\!\mathcal{E}_{n+1}$, $\mathcal{S}(e)\!\setminus\!\{f\}\!\neq\!\emptyset$ by Eq.~\eqref{eq:p2-orph}, giving (I1); and for every $h\!\in\!\mathcal{H}_{n+1}$, the second filter forces $\mathcal{E}_h\!\subseteq\!\mathcal{E}_n\!\setminus\!\mathcal{E}_n^{\mathrm{orph}}(f)\!=\!\mathcal{E}_{n+1}$, giving (I2).

When $\Delta_n\!=\!\textsc{Refresh}(h_{\mathrm{old}}\!\to\!h_{\mathrm{new}})$, Reason synthesises the replacement proposition from the conflict context and Eq.~\eqref{eq:refresh} re-embeds in place,
\begin{equation}
\begin{aligned}
&\pi_{h_{\mathrm{new}}}\;=\;\mathrm{LLM}_{\mathrm{Reason}}\!\bigl(\pi_{h_{\mathrm{old}}},\,\mathrm{Query}(\Delta_n)\bigr),\qquad \phi(h_{\mathrm{new}})\;=\;\mathrm{Embed}(\pi_{h_{\mathrm{new}}}),\\[2pt]
&\mathcal{E}_{h_{\mathrm{new}}}\!=\!\mathcal{E}_{h_{\mathrm{old}}}\!\subseteq\!\mathcal{E}_n,\qquad \mathcal{A}_{h_{\mathrm{new}}}\!=\!\mathcal{A}_{h_{\mathrm{old}}},\qquad \omega_{h_{\mathrm{new}}}\!\in\![0,1],\\[2pt]
&\mathcal{E}_{n+1}\!=\!\mathcal{E}_n,\qquad \mathcal{H}_{n+1}\!=\!(\mathcal{H}_n\!\setminus\!\{h_{\mathrm{old}}\})\!\cup\!\{h_{\mathrm{new}}\},\qquad \mathcal{A}_{n+1}\!=\!\mathcal{A}_n.
\end{aligned}
\label{eq:p2-ref-state}
\end{equation}
The first line preserves the embedding-coherence invariant since $\phi$ is recomputed from the new proposition; the second line preserves the vertex set $\mathcal{E}_{h_{\mathrm{new}}}$ and the support $\mathcal{A}_{h_{\mathrm{new}}}$ verbatim. (I1) holds vacuously because no entity is removed, and (I2) holds because $\mathcal{E}_{h_{\mathrm{new}}}\!\subseteq\!\mathcal{E}_n\!=\!\mathcal{E}_{n+1}$.

The three branches above exhaust the event alphabet, and Eqs.~\eqref{eq:p2-add-state}--\eqref{eq:p2-ref-state} each preserve (I1)--(I2). Hence by induction (I1) and (I2) hold at $\mathcal{G}_N$, so the audit log $\mathcal{L}_N\!=\!\bigcup_{n<N}\!\{(\Delta_n,\mathcal{G}_{n+1}\!\ominus\!\mathcal{G}_n)\}$ records a sequence of states each free of orphan entities and dangling hyperedges, establishing Proposition~2.
\end{proof}

\subsection{Proof of Proposition 3}\label{proof3}
\textbf{Proposition 3.} \textit{Multi-angle routing $\pi^{*}(q)$ covers $\{T3,T4,T6\}$, while no single $a\!\in\!\mathcal{R}$ does.}\vspace{-2mm}
\begin{proof}
Let $\mathcal{Q}_t$ denote the query population for task $t\!\in\!\{T3,T4,T6\}$ and $\mathrm{score}_t(R)\!\in\![0,1]$ the task-specific score on retrieval output $R$. We say angle $a\!\in\!\mathcal{R}\!=\!\{\mathrm{sub},\mathrm{com},\mathrm{disc}\}$ \emph{covers} $t$ when
\begin{equation}
\inf_{q\in\mathcal{Q}_t}\mathrm{score}_t\!\bigl(R_a(q)\bigr)\;>\;0.
\label{eq:p3-cover}
\end{equation}
We prove (i) each task admits a covering angle, (ii) no single angle covers all three, and (iii) $\pi^{*}$ covers them jointly.

For (i), take a 3-hop factual query $q_3\!\in\!\mathcal{Q}_{T3}$ with gold path $\gamma_{q_3}\!=\!(e_0,h_1,e_1,\dots,h_3,e_3)\!\subseteq\!\mathcal{G}$. The reciprocal-rank fusion of Eq.~\eqref{eq:sub} ensures that at least the path endpoints fall inside the seeds, and $\ell$-hop expansion then closes the path:
\begin{equation}
\begin{aligned}
\mathcal{Z}_{q_3}\;=\;\mathrm{Top}_{K}\!\Bigl[\,s_{\mathrm{RRF}}(\,\cdot\,;q_3)\,\Bigr]\;\supseteq\;\{e_0,e_3\},\qquad K\!\geq\!2,\\[2pt]
\mathcal{N}_{\leq\ell}(\mathcal{Z}_{q_3})\;=\;\bigcup_{x\in\mathcal{Z}_{q_3}}\bigcup_{i=0}^{\ell}\!\mathcal{N}^{(i)}(x)\;\supseteq\;\{e_0,h_1,e_1,h_2,e_2,h_3,e_3\}\;=\;\gamma_{q_3},\qquad \ell\!\geq\!2,\\[2pt]
R_{\mathrm{sub}}(q_3)\;=\;\mathcal{G}\!\bigl[\mathcal{N}_{\leq\ell}(\mathcal{Z}_{q_3})\bigr]\;\supseteq\;\mathcal{G}[\gamma_{q_3}],
\end{aligned}
\label{eq:p3-T3}
\end{equation}
so $R_{\mathrm{sub}}(q_3)$ contains the gold spans, yielding $\mathrm{score}_{T3}\!\bigl(R_{\mathrm{sub}}(q_3)\bigr)\!=\!\mathrm{F1}_{\mathrm{tok}}(R_{\mathrm{sub}}(q_3),\gamma_{q_3})\!>\!0$. For a cross-modal query $q_4\!\in\!\mathcal{Q}_{T4}$ on a modality subset $\mathcal{M}'$ with $|\mathcal{M}'|\!\geq\!2$, Proposition~1 implies that only $\mathcal{G}$ encodes $\mu(\mathcal{A}_h)$, and T4's score decomposes as
\begin{equation}
\mathrm{score}_{T4}(R)\;=\;0.7\!\cdot\!\mathrm{F1}(R)\;+\;0.3\!\cdot\!\mathrm{XMEv}(R),\qquad \mathrm{XMEv}(R)\!>\!0\;\Longleftrightarrow\;\exists\,h\!\in\!R\,:\,|\mu(\mathcal{A}_h)|\!\geq\!2,
\label{eq:p3-T4}
\end{equation}
so $R_{\mathrm{sub}}(q_4)$ scores positively whenever $\mathcal{G}_{q_4}^{\mathrm{sub}}$ contains a multi-modal hyperedge. For a topic-discovery query $q_6\!\in\!\mathcal{Q}_{T6}$, Eq.~\eqref{eq:disc-route} produces the gap and suggestion components exclusively through $\mathrm{LLM}_{\mathrm{disc}}$:
\begin{equation}
\bigl(\mathrm{gaps}(t),\mathrm{sug}(t)\bigr)\;=\;\mathrm{LLM}_{\mathrm{disc}}\!\bigl(t,\,\{\sigma_k\}_{\mathcal{C}_t^{\mathrm{com}}},\,\mathrm{known}(t)\bigr),
\label{eq:p3-T6}
\end{equation}
so the bigram-presence test on $\mathrm{gaps},\mathrm{sug}$ gives $\mathrm{score}_{T6}\!\bigl(R_{\mathrm{disc}}(q_6)\bigr)\!>\!0$. Thus $\mathrm{sub}$ covers $T3$ and $T4$, and $\mathrm{disc}$ covers $T6$.

For (ii), three explicit failures suffice. The output schemas of the three angles project onto disjoint sub-components of the per-task scoring functions, so each angle leaves at least one task with zero score:
\begin{equation}
\begin{aligned}
\mathrm{sub}\;\text{on}\;q_6\!:\;\;& R_{\mathrm{sub}}(q_6)\!\cap\!\bigl\{\mathrm{gaps}(t),\mathrm{sug}(t)\bigr\}\!=\!\emptyset
\;\Longrightarrow\; \mathrm{score}_{T6}\!\bigl(R_{\mathrm{sub}}(q_6)\bigr)\!=\!\mathrm{Bigram}(\emptyset,\mathrm{KB})\!=\!0,\\[2pt]
\mathrm{com}\;\text{on}\;q_3\!:\;\;& R_{\mathrm{com}}(q_3)\!=\!\bigoplus_{k}\sigma_k\;\text{averages over}\;\mathcal{C}_k\!\not\supseteq\!\gamma_{q_3}
\!\Longrightarrow \mathrm{score}_{T3}\!\bigl(R_{\mathrm{com}}(q_3)\bigr)\!=\!\mathrm{F1}_{\mathrm{tok}}(\sigma,\gamma_{q_3})\!=\!0,\\[2pt]
\mathrm{disc}\;\text{on}\;q_3\!:\;\;& R_{\mathrm{disc}}(q_3)\!\in\!\bigl(\mathrm{known},\mathrm{gaps},\mathrm{sug}\bigr)\!\not\subseteq\!\text{answer-space of }T3
\;\Longrightarrow\; \mathrm{score}_{T3}\!\bigl(R_{\mathrm{disc}}(q_3)\bigr)\!=\!0.
\end{aligned}
\label{eq:p3-fails}
\end{equation}
Equivalently, every angle has a witness task on which it fails, $\forall a\!\in\!\mathcal{R}\;\exists t\!\in\!\{T3,T4,T6\}\!:\!\mathrm{score}_t(R_a(q))\!=\!0$, which forces a strict subset relation on the intersection of per-angle covered task sets:
\begin{equation}
\bigcap_{a\in\mathcal{R}}\bigl\{t\,:\,a\text{ covers }t\bigr\}\;\subseteq\;\{T3,T4,T6\}\!\setminus\!\bigl\{T6,T3,T3\bigr\}\;=\;\{T4\}\;\subsetneq\;\{T3,T4,T6\}.
\label{eq:p3-no-single}
\end{equation}

For (iii), the router from Eq.~\eqref{eq:disc-route} selects
\begin{equation}
\pi^{*}(q)\;=\;\{a\!\in\!\mathcal{R}\,:\,w_a(q)\!\geq\!\theta_{\pi}\},\qquad \hat{a}(q)\;=\;\mathrm{LLM}_{\mathrm{gen}}\!\Bigl(q,\!\!\bigoplus_{a\in\pi^{*}(q)}\!R_a(q)\Bigr),
\label{eq:p3-router}
\end{equation}
where $\bigoplus$ preserves all evidence content. Writing $a^{*}(t)$ for any covering angle of task $t$ from (i), the empirical confusion matrix in \S\ref{5.6} guarantees per-task router calibration:
\begin{equation}
\Pr_{q\sim\mathcal{Q}_t}\!\bigl[a^{*}(t)\!\in\!\pi^{*}(q)\bigr]\;=\;\Pr_{q\sim\mathcal{Q}_t}\!\bigl[w_{a^{*}(t)}(q)\!\geq\!\theta_{\pi}\bigr]\;\geq\;1\!-\!\varepsilon_t,\qquad \varepsilon_t\!<\!1,\;\;t\!\in\!\{T3,T4,T6\},
\label{eq:p3-router-emp}
\end{equation}
where the slack $\varepsilon_t$ matches the off-diagonal mass of the routing confusion matrix. Since $\bigoplus$ is monotone under evidence inclusion (the gold evidence remains present once any covering angle enters the concatenation), the per-task score is bounded below by the score of $R_{a^{*}(t)}(q)$ alone:
\begin{equation}
\inf_{q\in\mathcal{Q}_t}\mathrm{score}_t\!\Bigl(\!\bigoplus_{a\in\pi^{*}(q)}\!R_a(q)\Bigr)\;\geq\;\inf_{q\in\mathcal{Q}_t}\mathrm{score}_t\!\bigl(R_{a^{*}(t)}(q)\bigr)\;>\;0,\qquad t\!\in\!\{T3,T4,T6\},
\label{eq:p3-joint}
\end{equation}
where $a^{*}(t)$ denotes any covering angle from (i).

Combining Eq.~\eqref{eq:p3-no-single} with Eq.~\eqref{eq:p3-joint}, the routed concatenation $\pi^{*}$ covers $\{T3,T4,T6\}$ jointly while no single $a\!\in\!\mathcal{R}$ does, establishing Proposition~3.
\end{proof}

\clearpage

\section{KnowledgeClaw Algorithm Detail}\label{app:algo}

This appendix supplements \S\ref{sec:method} with executable pseudocode for the three contributions, mirroring the algorithm-card style of HyperGraphRAG~\citep{HyperGraphRAG}. Algorithm~\ref{alg:kc-c1} instantiates Eq.~\eqref{eq:hg} by populating $\mathcal{G}$ from a heterogeneous corpus $\mathcal{A}$; Algorithm~\ref{alg:kc-c2} instantiates Eq.~\eqref{eq:loop} as a transactional \textsc{Build}--\textsc{Query}--\textsc{Reason}--\textsc{Update} loop whose \textsc{Del}/\textsc{Refresh} branches realise Eqs.~\eqref{eq:cascade}--\eqref{eq:refresh}; Algorithm~\ref{alg:kc-c3} instantiates Eqs.~\eqref{eq:sub}--\eqref{eq:disc-route} by routing every $q$ to a calibrated subset of $\mathcal{R}\!=\!\{\mathrm{sub},\mathrm{com},\mathrm{disc}\}$ before generation.

\textbf{C1: Omnimodal Hypergraph Construction.} Each asset $f\!\in\!\mathcal{A}$ is first lifted to a propositional view by the modality-aware extractor $\mathrm{Extract}_{\mathrm{LLM}}$, which emits the candidate entities, propositions, hyperedges, and observed-modality map. Within-modality canonicalization $\kappa$ then folds surface variants (abbreviations, paraphrases) of the same entity into one canonical node, after which every still-disjoint cross-modality pair $(e_a,e_b)$ is gated by the two-step test $\chi$ of Eq.~\eqref{eq:coref}: a cheap cosine pre-filter at threshold $\theta_{\chi}$ followed by an LLM verifier on the survivors. Accepted pairs become explicit cross-modal hyperedges, so $\mathcal{G}$ is persisted as a bipartite incidence matrix between $\mathcal{E}$ and $\mathcal{H}$ with shared dense embeddings $\phi$ that all later retrieval stages reuse.

\begin{algorithm}[h!t]
\small
\caption{C1 --- Omnimodal Knowledge Hypergraph Construction}
\label{alg:kc-c1}
\begin{algorithmic}[1]
\REQUIRE Asset corpus $\mathcal{A}\!\subseteq\!\mathcal{D}$, modality probe $m(\cdot)$, coreference threshold $\theta_{\chi}$
\ENSURE Typed omnimodal hypergraph $\mathcal{G}\!=\!(\mathcal{E},\mathcal{H},\mathcal{A},\mathcal{P};\mu,\tau,\chi,\phi)$
\STATE Initialize $\mathcal{E}\!\leftarrow\!\emptyset,\;\mathcal{H}\!\leftarrow\!\emptyset,\;\mathcal{P}\!\leftarrow\!\emptyset$ and modality map $\mu\!\leftarrow\!\emptyset$
\FOR{each asset $f\!\in\!\mathcal{A}$}
    \STATE $(\hat{\mathcal{E}}_f,\hat{\mathcal{P}}_f,\hat{\mathcal{H}}_f,\hat{\mu}_f)\!\leftarrow\!\mathrm{Extract}_{\mathrm{LLM}}(f,m(f))$ \hfill \COMMENT{T1: per-modality propositional extraction}
    \STATE Update $\mu(e)\!\leftarrow\!\mu(e)\!\cup\!\hat{\mu}_f(e)$ for every $e\!\in\!\hat{\mathcal{E}}_f$
    \STATE $\mathcal{E}\!\leftarrow\!\mathcal{E}\!\cup\!\kappa(\hat{\mathcal{E}}_f),\;\mathcal{P}\!\leftarrow\!\mathcal{P}\!\cup\!\hat{\mathcal{P}}_f,\;\mathcal{H}\!\leftarrow\!\mathcal{H}\!\cup\!\hat{\mathcal{H}}_f$
\ENDFOR
\STATE Compute embeddings $\phi(x)\!\leftarrow\!\mathrm{Embed}(x)$ for every $x\!\in\!\mathcal{E}\!\cup\!\mathcal{H}$
\FOR{each pair $(e_a,e_b)\!\in\!\mathcal{E}\!\times\!\mathcal{E}$ with $\mu(e_a)\!\cap\!\mu(e_b)\!=\!\emptyset$}
    \IF{$\langle\phi(e_a),\phi(e_b)\rangle\!\geq\!\theta_{\chi}$ \AND $\mathrm{Verify}_{\mathrm{LLM}}(e_a,e_b\,|\,\mathcal{S}(e_a)\!\cup\!\mathcal{S}(e_b))$}
        \STATE Insert cross-modal hyperedge $h_{ab}\!=\!(\{e_a,e_b\},\textit{coref},\pi_{ab},\mathcal{S}(e_a)\!\cup\!\mathcal{S}(e_b),1)$, set $\chi(e_a,e_b)\!\leftarrow\!1$ \hfill \COMMENT{T2}
        \STATE $\mathcal{H}\!\leftarrow\!\mathcal{H}\!\cup\!\{h_{ab}\}$
    \ENDIF
\ENDFOR
\STATE Persist $\mathcal{G}$ as bipartite incidence $M\!\in\!\{0,1\}^{|\mathcal{E}|\times|\mathcal{H}|}$ with vector indices over $\phi$
\RETURN $\mathcal{G}$
\end{algorithmic}
\end{algorithm}

\textit{Complexity Analysis.} With $|\mathcal{A}|\!=\!N$ assets, $r$ propositions per asset, and $n$ entities per proposition, extraction and embedding cost $\mathcal{O}(N\!\cdot\!r\!\cdot\!n)$ remote calls (parallelisable). The cross-modal pair scan is reduced from quadratic to $\mathcal{O}(|\mathcal{E}|\!\cdot\!\log|\mathcal{E}|)$ via an HNSW index that returns only the top-$\kappa$ neighbours above $\theta_{\chi}$, after which $\mathrm{Verify}_{\mathrm{LLM}}$ fires on the $\mathcal{O}(\kappa\!\cdot\!|\mathcal{E}|)$ survivors --- this is what keeps T2 tractable on the $\approx\!4{,}500$-asset OmniKE-Bench.

\textbf{C2: Closed-loop Lifecycle Agent.} Every corpus event $\Delta_n\!\in\!\{\textsc{Add},\textsc{Del},\textsc{Refresh}\}\!\times\!\mathcal{D}$ drives a single transactional pass through the four-stage operator of Eq.~\eqref{eq:loop}: \textsc{Build} ingests $f$ into a candidate fragment, \textsc{Query} pulls the affected neighbourhood out of the bipartite index, \textsc{Reason} asks the LLM to detect conflicts and synthesise the structured diff, and \textsc{Update} commits $\mathcal{G}_{n+1}$ together with the audit-log entry $(\Delta_n,\mathcal{G}_{n+1}\!\ominus\!\mathcal{G}_n)$ used for replay. The non-trivial branches realise Eqs.~\eqref{eq:cascade}--\eqref{eq:refresh}: \textsc{Del} performs cascading orphan eviction so no entity outlives its only support and no hyperedge is left dangling, while \textsc{Refresh} rewrites $h_{\mathrm{old}}\!\to\!h_{\mathrm{new}}$ in place and re-embeds via $\phi(h_{\mathrm{new}})\!=\!\mathrm{Embed}(\pi_{h_{\mathrm{new}}})$, tracking the new fact without disturbing neighbours.

\begin{algorithm}[h!t]
\small
\caption{C2 --- Closed-loop Lifecycle Agent (Build--Query--Reason--Update)}
\label{alg:kc-c2}
\begin{algorithmic}[1]
\REQUIRE Live hypergraph $\mathcal{G}_n$, event $\Delta_n\!=\!(\mathrm{op},f)$, audit log $\mathcal{L}_n$
\ENSURE Updated $\mathcal{G}_{n+1}$, log $\mathcal{L}_{n+1}$
\STATE \textbf{Build:} parse $f$, run modality probe $m(f)$, and produce candidate $(\hat{\mathcal{E}}_f,\hat{\mathcal{P}}_f,\hat{\mathcal{H}}_f)$ via $\mathrm{Extract}_{\mathrm{LLM}}$
\STATE \textbf{Query:} retrieve neighborhood $\mathcal{N}(f)\!=\!\{h\!\in\!\mathcal{H}_n\!:\,\mathcal{A}_h\!\ni\!f\!\vee\!\mathcal{E}_h\!\cap\!\hat{\mathcal{E}}_f\!\neq\!\emptyset\}$ via the bipartite index
\STATE \textbf{Reason:} call $\mathrm{LLM}_{\mathrm{Reason}}(\hat{\mathcal{H}}_f,\mathcal{N}(f))$ to detect conflicts and emit a structured diff $\delta_n$
\IF{$\mathrm{op}\!=\!\textsc{Add}$}
    \STATE Insert $\hat{\mathcal{E}}_f$ via $\kappa$, append $\hat{\mathcal{H}}_f$, and re-test cross-modal $\chi$ on $\hat{\mathcal{E}}_f\!\times\!\mathcal{E}_n$
\ELSIF{$\mathrm{op}\!=\!\textsc{Del}$}
    \STATE Compute orphans $\mathcal{E}_n^{\mathrm{orph}}(f)\!=\!\{e\!\in\!\mathcal{E}_n\!:\,\mathcal{S}(e)\!\subseteq\!\{f\}\}$ \hfill \COMMENT{single-support entities}
    \STATE $\mathcal{E}_{n+1}\!\leftarrow\!\mathcal{E}_n\!\setminus\!\mathcal{E}_n^{\mathrm{orph}}(f)$
    \STATE $\mathcal{H}_{n+1}\!\leftarrow\!\{h\!\in\!\mathcal{H}_n\!:\,\mathcal{A}_h\!\not\subseteq\!\{f\}\,\wedge\,\mathcal{E}_h\!\cap\!\mathcal{E}_n^{\mathrm{orph}}(f)\!=\!\emptyset\}$ \hfill \COMMENT{drop dangling hyperedges}
    \STATE $\mathcal{A}_{n+1}\!\leftarrow\!\mathcal{A}_n\!\setminus\!\{f\}$
\ELSIF{$\mathrm{op}\!=\!\textsc{Refresh}(h_{\mathrm{old}}\!\to\!h_{\mathrm{new}})$}
    \STATE $\pi_{h_{\mathrm{new}}}\!\leftarrow\!\mathrm{LLM}_{\mathrm{Reason}}(\pi_{h_{\mathrm{old}}},\mathcal{N}(f))$, $\phi(h_{\mathrm{new}})\!\leftarrow\!\mathrm{Embed}(\pi_{h_{\mathrm{new}}})$
    \STATE $\mathcal{H}_{n+1}\!\leftarrow\!(\mathcal{H}_n\!\setminus\!\{h_{\mathrm{old}}\})\!\cup\!\{h_{\mathrm{new}}\}$
\ENDIF
\STATE \textbf{Update:} commit transactionally, log diff $\mathcal{L}_{n+1}\!\leftarrow\!\mathcal{L}_n\!\cup\!\{(\Delta_n,\mathcal{G}_{n+1}\!\ominus\!\mathcal{G}_n)\}$
\RETURN $(\mathcal{G}_{n+1},\mathcal{L}_{n+1})$
\end{algorithmic}
\end{algorithm}

\textit{Complexity Analysis.} Each event touches only the local neighbourhood of $f$: \textsc{Build}, \textsc{Query}, \textsc{Reason}, \textsc{Update} cost one extractor call, $\mathcal{O}(\deg(f))$ over the bipartite index, $\mathcal{O}(|\hat{\mathcal{H}}_f|\!+\!|\mathcal{N}(f)|)$, and $\mathcal{O}(|\mathcal{E}_n^{\mathrm{orph}}(f)|\!+\!|\hat{\mathcal{H}}_f|)$ respectively. The full pass is therefore amortised $\mathcal{O}(\deg(f))$ rather than $\mathcal{O}(|\mathcal{G}_n|)$, so serving latency stays stable as the graph grows; with Proposition~2 this produces the flat T5(t) curve in \S\ref{5.5} and forbids any orphan or dangling hyperedge from leaking between events.

\textbf{C3: Multi-angle Application with Intent Routing.} The serving path scores all three angles with one router pass on $q$, materialises only those whose score $w_a(q)$ clears the gate $\theta_{\pi}$, and concatenates their outputs into the single context window seen by the generator. \emph{Subgraph retrieval} fuses entity- and hyperedge-rank lists by reciprocal-rank fusion and grows the seed set by an $\ell$-hop walk over the bipartite index (Eq.~\eqref{eq:sub}); \emph{community synthesis} reuses the Louvain partition cached at build time and ranks summaries against $\phi(q)$ (Eq.~\eqref{eq:com}); \emph{knowledge discovery} produces the $(\mathrm{known},\mathrm{gaps},\mathrm{sug})$ triple required by T6 (Eq.~\eqref{eq:disc-route}). By Proposition~3 the calibrated $w_a(q)$ fires only one or two angles per query on average, not all three.

\begin{algorithm}[h!t]
\small
\caption{C3 --- Intent Routing and Multi-angle Retrieval-Generation}
\label{alg:kc-c3}
\begin{algorithmic}[1]
\REQUIRE Query $q$, hypergraph $\mathcal{G}$, communities $\{(\mathcal{C}_k,\sigma_k)\}$, RRF constant $k_0$, walk depth $\ell$, router gate $\theta_{\pi}$, discovery threshold $\theta_{\mathrm{k}}$
\ENSURE Final answer $\hat{a}(q)$
\STATE Compute angle scores $w_a(q)\!\leftarrow\!\mathrm{Router}(q,a)$ for every $a\!\in\!\mathcal{R}\!=\!\{\mathrm{sub},\mathrm{com},\mathrm{disc}\}$
\STATE Select angles $\pi^{*}(q)\!\leftarrow\!\{a\!\in\!\mathcal{R}\!:\,w_a(q)\!\geq\!\theta_{\pi}\}$; if $\pi^{*}(q)\!=\!\emptyset$ fall back to $\arg\max_a w_a(q)$
\IF{$\mathrm{sub}\!\in\!\pi^{*}(q)$}
    \STATE $r_{\mathcal{E}}(\cdot;q),r_{\mathcal{H}}(\cdot;q)\!\leftarrow\!$ cosine ranks of $\phi(q)$ over $\phi(\mathcal{E}),\phi(\mathcal{H})$
    \STATE $s_{\mathrm{RRF}}(x;q)\!\leftarrow\!\sum_{r}\!1/(k_0\!+\!r(x;q))$, take top-$K$ seeds $\mathcal{Z}_q$
    \STATE Expand $\mathcal{N}_{\leq\ell}(\mathcal{Z}_q)\!\leftarrow\!\bigcup_{x\in\mathcal{Z}_q}\!\bigcup_{i\leq\ell}\!\mathcal{N}^{(i)}(x)$ and set $R_{\mathrm{sub}}\!\leftarrow\!\mathcal{G}[\mathcal{N}_{\leq\ell}(\mathcal{Z}_q)]$ \hfill \COMMENT{T3, T4}
\ENDIF
\IF{$\mathrm{com}\!\in\!\pi^{*}(q)$}
    \STATE $R_{\mathrm{com}}\!\leftarrow\!\mathrm{TopL}_{\mathcal{C}_k}\!\bigl[\langle\phi(q),\phi(\sigma_k)\rangle\bigr]$ over the cached Louvain partition of $A_{ij}\!=\!\sum_{h:\{i,j\}\subseteq\mathcal{E}_h}\!\omega_h$
\ENDIF
\IF{$\mathrm{disc}\!\in\!\pi^{*}(q)$}
    \STATE $\mathrm{known}(q)\!\leftarrow\!\{e\!\in\!\mathcal{E}\!:\,\langle\phi(q),\phi(e)\rangle\!\geq\!\theta_{\mathrm{k}}\}$
    \STATE $(\mathrm{gaps},\mathrm{sug})\!\leftarrow\!\mathrm{LLM}_{\mathrm{disc}}(q,\{\sigma_k\}_{\mathcal{C}_q^{\mathrm{com}}},\mathrm{known}(q))$, set $R_{\mathrm{disc}}\!\leftarrow\!(\mathrm{known},\mathrm{gaps},\mathrm{sug})$ \hfill \COMMENT{T6}
\ENDIF
\STATE Concatenate $R_{\pi^{*}}\!\leftarrow\!\bigoplus_{a\in\pi^{*}(q)}\!R_a$ preserving evidence order
\STATE $\hat{a}(q)\!\leftarrow\!\mathrm{LLM}_{\mathrm{gen}}(q,R_{\pi^{*}})$
\RETURN $\hat{a}(q)$
\end{algorithmic}
\end{algorithm}

\textit{Complexity Analysis.} Routing is a single short LLM call; subgraph retrieval is dominated by ANN lookup over $\phi(\mathcal{E})\!\cup\!\phi(\mathcal{H})$ at $\mathcal{O}(\log(|\mathcal{E}|\!+\!|\mathcal{H}|))$ per index plus an $\ell$-hop walk bounded by $\mathcal{O}(K\!\cdot\!\bar{d}^{\ell})$ on the average degree $\bar{d}$; community synthesis reuses the cached Louvain partition, costing $\mathcal{O}(|\{\mathcal{C}_k\}|)$ per query; and discovery costs one $\mathrm{LLM}_{\mathrm{disc}}$ call. Generation is treated as black-box $\mathcal{O}(L)$ in prompt length, which is what dominates the 580\,ms end-to-end serving latency reported in \S\ref{5.7}.

\textbf{End-to-end pipeline.} Algorithms~\ref{alg:kc-c1}--\ref{alg:kc-c3} compose as $\mathrm{C1}\!\to\!\mathrm{C2}^{\,*}\!\to\!\mathrm{C3}$: a one-time C1 pass yields $\mathcal{G}_0$, C2 then runs once per event $\Delta_n$ producing $(\mathcal{G}_{n+1},\mathcal{L}_{n+1})$, and serving issues C3 against the latest $\mathcal{G}_n$ without blocking on C2 since the bipartite index is copy-on-write. The three are thus both ablation-orthogonal (Table~\ref{tab:main}) and pipeline-orthogonal --- each can be cached or scaled independently --- yielding the cost--latency profile in \S\ref{5.7} (\$1.95/persona, 41\,min build, 580\,ms serving). With the propositions of Appendix~\ref{proofs}, Algorithms~\ref{alg:kc-c1}--\ref{alg:kc-c3} fully specify KnowledgeClaw's executable surface.
\clearpage

\section{OmniKE-Bench Detail}\label{app:bench}

This appendix expands \S\ref{sec:task} into a self-contained reference for OmniKE-Bench, the first end-to-end benchmark for agentic knowledge engineering on real omnimodal personal files. We document the four design axes in turn: task definitions (\S\ref{app:bench-task}), the construction pipeline that turns raw files into evaluable units (\S\ref{app:bench-pipe}), per-persona and per-modality corpus statistics (\S\ref{app:bench-stat}), and the per-task metric formulas (\S\ref{app:bench-metric}). Throughout, $\mathcal{M}\!=\!\{\textit{text},\textit{pdf},\textit{image},\textit{audio},\textit{video}\}$, $\mathcal{D}\!=\!\{f_i\}$ with $m(f_i)\!\in\!\mathcal{M}$, and the ten persona codes \textsc{HC}/\textsc{Fin}/\textsc{Edu}/\textsc{Law}/\textsc{Psy}/\textsc{Mus}/\textsc{Gam}/\textsc{Bld}/\textsc{Sft}/\textsc{Rob} match Table~\ref{tab:main}.

\subsection{Task Design}\label{app:bench-task}

The six tasks split into three groups, one per property of \S\ref{sec:intro}: \emph{Construction} (T1, T2) tests omnimodality, \emph{Maintenance} (T5) tests closed-loop lifecycle, and \emph{Application} (T3, T4, T6) tests multi-angle application. All six share the same corpus $\mathcal{D}$ and persona granularity. Each task is specified as a deterministic function $\mathcal{T}_j$ that a system under evaluation must implement, so a single OKE score on a persona reflects every property simultaneously rather than a paradigm-favoured subset.

\noindent\textbf{T1 (Knowledge Acquisition).} Lift the corpus to a propositional view:
\begin{equation}
\mathcal{T}_1\!:\,\mathcal{D}\;\longmapsto\;(\hat{\mathcal{E}},\hat{\mathcal{P}},\hat{\mu}),\qquad \hat{\mu}\!:\,\hat{\mathcal{E}}\!\to\!2^{\mathcal{M}},
\label{eq:t-t1}
\end{equation}
where $\hat{\mathcal{E}}$ is the extracted entity set, $\hat{\mathcal{P}}$ the proposition set, and $\hat{\mu}$ the per-entity modality map (an entity may be observed in several modalities). Reference-free; tests whether all five modalities are unified into a single symbolic substrate rather than collapsed onto text.

\noindent\textbf{T2 (Cross-modal Coreference).} Link mentions of the same referent across disjoint modalities:
\begin{equation}
\mathcal{T}_2\!:\,(\mathcal{D},\hat{\mathcal{E}})\;\longmapsto\;\hat{L}\;\subseteq\;\bigl\{(e_a,e_b)\!\in\!\hat{\mathcal{E}}^{2}\!:\,\hat{\mu}(e_a)\!\cap\!\hat{\mu}(e_b)\!=\!\emptyset\bigr\}.
\label{eq:t-t2}
\end{equation}
The disjointness constraint rules out within-modality string matches, so $\hat{L}$ is forced to express genuinely cross-modal identity (e.g., a name in PDF linked to its face in image). Reference-free.

\noindent\textbf{T3 (Multi-hop Retrieval).} Compose evidence across multiple linked facts:
\begin{equation}
\mathcal{T}_3\!:\,(q,\mathcal{D})\;\longmapsto\;(\hat{a},R),\qquad q\!\in\!\mathcal{Q}_{2\text{--}4},\;\;R\!\subseteq\!\mathcal{D},
\label{eq:t-t3}
\end{equation}
where $\mathcal{Q}_{2\text{--}4}$ is the family of questions whose gold answer requires composing between $2$ and $4$ propositions, with each $q$ paired with a gold answer $a^{*}$ and the supporting evidence chain.

\noindent\textbf{T4 (Cross-modal Reasoning).} A modality-restricted variant of T3 that additionally exposes the cited modalities:
\begin{equation}
\mathcal{T}_4\!:\,(q,\mathcal{D})\;\longmapsto\;(\hat{a},R,\hat{\mu}(R)),\qquad \mu^{*}(q)\!\subseteq\!\mathcal{M},\;\;|\mu^{*}(q)|\!\geq\!2,
\label{eq:t-t4}
\end{equation}
with $\mu^{*}(q)$ the modalities the gold answer requires and $\hat{\mu}(R)\!=\!\bigcup_{f\in R}m(f)$ the modalities actually cited; a system returning the right answer from a single modality is penalised.

\noindent\textbf{T5 (Lifecycle Consistency).} A two-stage probe over the lifecycle operator:
\begin{equation}
\mathcal{T}_5\!:\,(\mathcal{G}_n,\Delta_n,q)\;\longmapsto\;\hat{a},\qquad \Delta_n\!\in\!\{\textsc{Add},\textsc{Del},\textsc{Refresh}\}\!\times\!\mathcal{D},
\label{eq:t-t5}
\end{equation}
i.e., the system first applies $\Delta_n$ to its live state and then answers $q$, whose correct response depends on whether $\Delta_n$ was committed. The two probe families are \emph{NL-CRUD} (free-form instructions plus a binary verification query) and \emph{Fact-refresh} (the gold answer must shift from $h_{\mathrm{old}}$ to $h_{\mathrm{new}}$).

\noindent\textbf{T6 (Knowledge Discovery).} Emit a discovery triple over a topic:
\begin{equation}
\mathcal{T}_6\!:\,t\;\longmapsto\;\bigl(\mathrm{known}(t),\,\mathrm{gaps}(t),\,\mathrm{sug}(t)\bigr),\quad \mathrm{known}(t)\!\subseteq\!\hat{\mathcal{E}},\;\;\mathrm{gaps}(t),\,\mathrm{sug}(t)\!\not\subseteq\!\hat{\mathcal{E}}.
\label{eq:t-t6}
\end{equation}
$\mathrm{known}(t)$ enumerates entities already covered, $\mathrm{gaps}(t)$ entities that should plausibly be in $\mathcal{D}$ but are missing, and $\mathrm{sug}(t)$ exploratory directions outside $\mathcal{D}$. Reference-free; tests whether the application layer can perform proactive inference over communities rather than passive lookup against the index.

\subsection{Benchmark Construction}\label{app:bench-pipe}

OmniKE-Bench is built in four phases, executed once per persona; the entire pipeline is reproducible from the released scripts. Total budget per persona: $\approx$3\,min file collection (manual), $\approx$30\,min derivation (parallel API calls), $\approx$\$3 annotation, $<$1\,min freezing.

\noindent\textbf{(1) Corpus collection.} Each persona is a researcher profile populated by harvesting publicly available files matching the domain (e.g., \textsc{HC} = clinical guidelines + medical-image atlases + grand-rounds recordings + procedure videos). Files are kept in their native formats, never pre-converted to text. Across the ten personas, the corpus contains \textbf{4{,}394} files totalling $\approx$22\,GB (985 text, 1{,}285 PDF, 1{,}507 image, 362 audio, 255 video; per-persona breakdown in Table~\ref{tab:bench-stats}). Modality balance is intentionally domain-skewed: \textsc{Mus}/\textsc{Rob} are audio/video-heavy, \textsc{Bld}/\textsc{Gam} image-heavy, \textsc{Law}/\textsc{HC} PDF-heavy, exposing whether a system absorbs low-frequency modalities rather than collapsing them onto text.

\noindent\textbf{(2) Derive.} Every $f$ is converted to a unified JSON record $\rho(f)$ with seven fields ---  \texttt{id}, \texttt{modality}, \texttt{path}, \texttt{text}, \texttt{caption}, \texttt{transcript}, \texttt{meta} --- via modality-specific extractors. Plain text is read raw; PDFs are parsed by \texttt{PyMuPDF} (text plus figure bounding boxes); images are captioned by \texttt{gpt-4o-mini} vision at temperature $0$ with a 256-token cap; audio is downsampled to 16\,kHz mono via \texttt{ffmpeg} and transcribed by \texttt{gemini-3.1-flash-lite-preview}; video is split into a 16\,kHz audio track plus key-frame captions sampled at 1\,fps. The resulting JSON forms the single substrate that every baseline and KnowledgeClaw later consume, so modality-handling differences cannot become a confound between methods evaluated on otherwise identical input.

\noindent\textbf{(3) Annotate.} The reference-free tasks (T1, T2, T6) score against $\rho(\mathcal{D})$ directly, so no gold is generated. The remaining three tasks are gold-labeled by \texttt{gpt-4o-mini}. T3 produces 50 multi-hop questions per persona by sampling 2--4-entity chains from a draft hypergraph and prompting the LLM to phrase a question whose answer is the chain's terminal entity. T4 produces 30 cross-modal questions per persona by restricting the chain to span $\geq\!2$ modalities and emitting both the answer and the required modality set $\mu^{*}(q)$. T5 produces 30 NL-CRUD probes (delete\,/\,update\,/\,other in a $3\!:\!2\!:\!1$ ratio) plus 20 fact-refresh tuples. T6 emits $\sim$35 topic descriptors per persona, drawn from the persona's high-degree entities. Every gold item is sanity-checked against $\rho(\mathcal{D})$ by an automatic round-trip: the answer must be locatable in the corpus, and the required modalities must actually appear in at least one of the cited evidence files.

\noindent\textbf{(4) Freeze.} Each persona is packaged as a self-contained directory containing the raw corpus, the derived JSON, the gold labels, and an evaluation config, and released alongside a deterministic CPU-only evaluator. The freeze step further records the SHA-256 of every file in the package so that future re-runs verify byte-equality of the entire input across deployment machines.

\subsection{Corpus Statistics}\label{app:bench-stat}

Table~\ref{tab:bench-stats} reports per-persona, per-modality file counts (left block) and per-task gold-item counts (right block). The corpus is intentionally heterogeneous: image is the dominant modality in absolute terms (1{,}507 files), driven by \textsc{Bld}/\textsc{Gam}/\textsc{Sft}; audio/video concentrate in \textsc{Mus}/\textsc{Rob}; \textsc{Law}/\textsc{HC} are PDF-heavy. Gold counts are kept uniform across personas (50/30/50/35 for T3/T4/T5/T6) so persona size does not inflate the OKE score.

\begin{table}[h!t]
\centering
\caption{OmniKE-Bench corpus and gold-label statistics. Left block: file counts per modality per persona; right block: gold-item counts per task per persona. T1 and T2 are reference-free and therefore have no gold items.}
\label{tab:bench-stats}
\setlength{\tabcolsep}{4.5pt}\renewcommand{\arraystretch}{1.05}
\small
\begin{tabular}{l|rrrrr|r||rrrr|r}
\toprule
& \multicolumn{6}{c||}{\textbf{Files by modality}} & \multicolumn{5}{c}{\textbf{Gold items by task}} \\
\textbf{Persona} & \textbf{text} & \textbf{pdf} & \textbf{image} & \textbf{audio} & \textbf{video} & \textbf{Sub.} & \textbf{T3} & \textbf{T4} & \textbf{T5} & \textbf{T6} & \textbf{Sub.} \\
\midrule
\textsc{HC}  & 95  & 175    & 100    & 30  & 18  & 418   & 50  & 30  & 50  & 35  & 165 \\
\textsc{Fin} & 110 & 160    & 130    & 20  & 22  & 442   & 50  & 30  & 50  & 35  & 165 \\
\textsc{Edu} & 105 & 130    & 145    & 30  & 22  & 432   & 50  & 30  & 50  & 35  & 165 \\
\textsc{Law} & 100 & 220    & 90     & 25  & 12  & 447   & 50  & 30  & 50  & 35  & 165 \\
\textsc{Psy} & 100 & 145    & 135    & 35  & 20  & 435   & 50  & 30  & 50  & 35  & 165 \\
\textsc{Mus} & 80  & 80     & 130    & 100 & 45  & 435   & 50  & 30  & 50  & 35  & 165 \\
\textsc{Gam} & 95  & 75     & 200    & 22  & 50  & 442   & 50  & 30  & 50  & 35  & 165 \\
\textsc{Bld} & 90  & 105    & 230    & 12  & 10  & 447   & 50  & 30  & 50  & 35  & 165 \\
\textsc{Sft} & 130 & 110    & 175    & 22  & 12  & 449   & 50  & 30  & 50  & 35  & 165 \\
\textsc{Rob} & 80  & 85     & 172    & 66  & 44  & 447   & 50  & 30  & 50  & 35  & 165 \\
\midrule
\textbf{Total} & \textbf{985} & \textbf{1{,}285} & \textbf{1{,}507} & \textbf{362} & \textbf{255} & \textbf{4{,}394} & \textbf{500} & \textbf{300} & \textbf{500} & \textbf{350} & \textbf{1{,}650} \\
\bottomrule
\end{tabular}
\end{table}

\subsection{Evaluation Metric Design}\label{app:bench-metric}

All six tasks are scored \emph{automatically} in $<\!1$\,s per persona without LLM calls, following the calibration protocol of RAGAS~\citep{RAGAS}; the reference-free tasks (T1, T2, T6) score against $\rho(\mathcal{D})$ rather than gold labels. Each $\mathrm{T}_j$ returns a value in $[0,100]$, and the aggregate \textbf{OKE} is the arithmetic mean over T1--T6 with architecturally unsupported tasks counted as $0$. We write $\rho(\mathcal{D})$ for the unified JSON substrate of \S\ref{app:bench-pipe}, $\mathrm{tok}(\cdot)$ for whitespace-and-punctuation tokenisation, and $\mathrm{Bigram}_{\geq3}(x,y)$ for the predicate that any $\geq\!3$-character word bigram of $x$ also occurs in $y$.

\noindent\textbf{T1 score.} Mean of four sub-scores in $[0,1]$:
\begin{equation}
\mathrm{T1} \;=\; 100\cdot\tfrac{1}{4}\bigl[\,F_{\mathcal{E}}\!+\!F_{\mathcal{P}}\!+\!Q_{\mathcal{E}}\!+\!C_{\mu}\,\bigr],
\label{eq:m-t1}
\end{equation}
with entity faithfulness $F_{\mathcal{E}}\!=\!\frac{1}{|\hat{\mathcal{E}}|}\!\sum_{e}\!\mathbb{1}[\mathrm{name}(e)\!\subseteq\!\rho(\mathcal{D})\!\vee\!\mathrm{TokOL}(e,\rho(\mathcal{D}))\!>\!0.2]$ blocking hallucinated names, proposition faithfulness $F_{\mathcal{P}}$ defined analogously over $\hat{\mathcal{P}}$, quantity sanity $Q_{\mathcal{E}}\!=\!\min(|\hat{\mathcal{E}}|/(\beta_{\mathcal{E}}|\mathcal{D}|),1)$ at target density $\beta_{\mathcal{E}}\!=\!8$ penalising empty extraction without rewarding inflation, and modality coverage $C_{\mu}\!=\!|\!\bigcup_e\!\hat{\mu}(e)|/|\mathcal{M}|$ rewarding reach across all five modalities.

\noindent\textbf{T2 score.} Mean of five binary checks, four averaged over $\hat{L}$ and one global:
\begin{equation}
\mathrm{T2} \;=\; 100\cdot\tfrac{1}{5}\bigl[\,F_L\!+\!R_L\!+\!N_L\!+\!D_L\!+\!Q_L\,\bigr],
\label{eq:m-t2}
\end{equation}
with $F_L$ (link faithfulness) requiring $\mathrm{Bigram}_{\geq3}(\mathrm{name}(e_a),\rho(\mathcal{D}))\!\wedge\!\mathrm{Bigram}_{\geq3}(\mathrm{name}(e_b),\rho(\mathcal{D}))$, $R_L$ (semantic relatedness) requiring $\langle\phi(e_a),\phi(e_b)\rangle\!\geq\!\theta_R\!=\!0.5$, $N_L$ (non-triviality) requiring $\mathrm{name}(e_a)\!\neq\!\mathrm{name}(e_b)$, $D_L$ (modality diversity) requiring $\hat{\mu}(e_a)\!\neq\!\hat{\mu}(e_b)$, and global $Q_L\!=\!\min(|\hat{L}|/\beta_L,1)$ at $\beta_L\!=\!|\hat{\mathcal{E}}|/4$ for link-quantity sanity, which jointly penalise both empty link sets and over-eager linkers.

\noindent\textbf{T3 score.} SQuAD-style token-level F1 between $\hat{a}$ and gold $a^{*}$:
\begin{equation}
\mathrm{T3} \;=\; 100\cdot F_1(\hat{a},a^{*}),\qquad F_1(\hat{a},a^{*})\!=\!\tfrac{2\,|\mathrm{tok}(\hat{a})\,\cap\,\mathrm{tok}(a^{*})|}{|\mathrm{tok}(\hat{a})|+|\mathrm{tok}(a^{*})|},
\label{eq:m-t3}
\end{equation}
matched after lower-casing and stop-word removal so that paraphrases sharing a content-word core score positively while answers that name the wrong entity collapse to $0$ regardless of how fluent they read.

\noindent\textbf{T4 score.} Convex combination of token-F1 and the cross-modal evidence rate:
\begin{equation}
\mathrm{T4} \;=\; 100\cdot\bigl[\,0.7\,F_1(\hat{a},a^{*})\,+\,0.3\,\mathrm{XMEv}(R)\,\bigr],\qquad \mathrm{XMEv}(R)\!=\!\frac{|\mu(R)\,\cap\,\mu^{*}(q)|}{|\mu^{*}(q)|},
\label{eq:m-t4}
\end{equation}
where the $0.7\!:\!0.3$ split keeps correctness dominant while requiring the cited modalities to cover $\mu^{*}(q)$; a single-modality answer to a $|\mu^{*}|\!=\!2$ question caps at $85$.

\noindent\textbf{T5 score.} Average of the two probe families:
\begin{equation}
\mathrm{T5} \;=\; 100\cdot\tfrac{1}{2}\bigl[\,\mathrm{NLC}_{\mathrm{acc}}\,+\,F_1(\hat{a}_{\mathrm{post}},a^{*}_{\mathrm{new}})\,\bigr],
\label{eq:m-t5}
\end{equation}
with $\mathrm{NLC}_{\mathrm{acc}}\!=\!\frac{1}{Z}\!\sum_{p}\!w_{\mathrm{op}(p)}\!\cdot\!\mathbb{1}[\mathrm{KW}(\mathrm{Resp}(p))]$ under weights $w_{\textsc{Del}}\!:\!w_{\textsc{Upd}}\!:\!w_{\mathrm{Oth}}\!=\!3\!:\!2\!:\!1$ ($Z\!=\!\sum_p w_{\mathrm{op}(p)}$, since deletes and updates are harder to commit consistently than additions), and the second term scoring the post-refresh answer against $a^{*}_{\mathrm{new}}$, so a CRUD path that commits the structural edit yet fails to propagate the new fact is still penalised.

\noindent\textbf{T6 score.} Bigram-presence test on the discovery triple:
\begin{equation}
\mathrm{T6} \;=\; 100\cdot\tfrac{1}{2}\bigl[\,\mathrm{Hit}_{\mathrm{known}}\,+\,(1\!-\!\mathrm{Hit}_{\mathrm{gaps}})\,\bigr],\qquad \mathrm{Hit}_X\!=\!\tfrac{1}{|X|}\!\sum_{x\in X}\!\mathbb{1}\!\bigl[\mathrm{Bigram}_{\geq3}(x,\rho(\mathcal{D}))\bigr],
\label{eq:m-t6}
\end{equation}
so $\mathrm{known}(t)$ is rewarded for items the corpus actually mentions while $\mathrm{gaps}(t)$ is rewarded for items it does not, separating proactive inference from corpus regurgitation; $\mathrm{sug}(t)$ is logged for qualitative inspection but not scored, since it intentionally lies outside $\mathcal{D}$.

\noindent\textbf{Aggregate.} The OKE score is $\mathrm{OKE}\!=\!\tfrac{1}{6}\!\sum_{j=1}^{6}\mathrm{T}_j$, with $\mathrm{T}_j\!=\!0$ when the system architecturally cannot produce the required output (e.g., a text-only RAG on T2/T4). This zero-on-unsupported convention is what penalises partial coverage of properties (i)--(iii) without separately re-tuning each baseline, and is the single rule that explains the large gaps between paradigm groups in Table~\ref{tab:main}.
\clearpage

\section{Baseline Detail}\label{app:baseline}

This appendix expands Table~\ref{tab:design} into a per-method specification. For each of the 22 baselines, the three KnowledgeClaw ablations, and KnowledgeClaw itself, we provide three blocks: \emph{(i)~Overview} explaining the method's original purpose and the substrate it materialises on $\rho(\mathcal{D})$; \emph{(ii)~Per-task implementation} formalising every $\mathcal{T}_j$ of \S\ref{app:bench-task} as an explicit operator, or annotating it as architecturally unsupported with a precise reason; and \emph{(iii)~Hyperparameters} listing the values used in our experiments. Methods are grouped by the seven paradigm classes of \S\ref{sec:exp} (\S\ref{app:b-chunk}--\S\ref{app:b-ake}), followed by the three KnowledgeClaw ablations (\S\ref{app:b-abl}) and KnowledgeClaw itself (\S\ref{app:b-kc}). Throughout we reuse $\rho(\mathcal{D})$ (the unified JSON substrate of \S\ref{app:bench-pipe}), $\phi(\cdot)$ (shared embeddings), $\mathrm{Chunk}(\cdot)$ (a fixed-size $300$-token chunker), $\mathrm{LLM}(\cdot)$ (the unified backbone \texttt{gemini-3.1-flash-lite-preview}), and $\mathrm{TopK}_{k}\!\langle\cdot,\cdot\rangle$ (top-$k$ retrieval under cosine). Whenever a method returns the empty set on $\mathcal{T}_j$ or fails its architectural prerequisite, the zero-on-unsupported convention of \S\ref{app:bench-metric} sets $\mathrm{T}_j\!=\!0$, exactly matching the ``-'' cells of Table~\ref{tab:design}. Each method below adopts the same default LLM backbone, embedding model, and zero-shot prompt as KnowledgeClaw to isolate architectural differences from prompting differences.

\subsection{Chunk/Hierarchical Retrieval}\label{app:b-chunk}

This family keeps the corpus as a flat or hierarchical chunk index without a symbolic entity layer, so only $\mathcal{T}_3$ (and partially $\mathcal{T}_6$ for tree-summary methods) is naturally supported.

\noindent\textbf{NaiveRAG~\citep{RAG}.}\quad\emph{(i)~Overview.} The canonical chunk-based RAG pipeline: each document is split into fixed-size token windows, every window is encoded by a sentence embedder, and the query retrieves the top-$k$ windows by cosine similarity for downstream LLM generation. Substrate: $\mathcal{I}\!=\!\{(c_i,\phi(c_i))\}_{c_i\in\mathrm{Chunk}(\mathcal{D})}$ (flat cosine index over the \texttt{text}/\texttt{caption}/\texttt{transcript} fields of $\rho(\mathcal{D})$).\quad\emph{(ii)~Per-task implementation.} $\mathcal{T}_1$: $(\hat{\mathcal{E}},\hat{\mathcal{P}},\hat{\mu})\!=\!(\emptyset,\emptyset,\emptyset)$ since chunks expose no symbolic entity layer; T1 receives only quantity-sanity and modality-coverage credit (Eq.~\eqref{eq:m-t1}). $\mathcal{T}_3$: $\hat{a}\!=\!\mathrm{LLM}\bigl(q,\,\mathrm{TopK}_{k}\!\langle\phi(q),\phi(c_i)\rangle\bigr)$. \emph{Unsupported}: $\mathcal{T}_2$ — $\hat{\mathcal{E}}\!=\!\emptyset$ admits no candidate pair $(e_a,e_b)$. $\mathcal{T}_4$ — every $c_i$ has $\mu(c_i)\!=\!\{\textit{text}\}$ after derivation, so $\mathrm{XMEv}(R)\!=\!0$ regardless of $\mu^{*}(q)$. $\mathcal{T}_5$ — the cosine index admits no $\Delta_n$ operator and is rebuilt from scratch on any change. $\mathcal{T}_6$ — no clusters or summaries, so $\mathrm{known}/\mathrm{gaps}$ have no anchor.\quad\emph{(iii)~Hyperparameters.} chunk size $300$ tokens, overlap $50$, retrieval top-$k\!=\!10$, no reranker; embedder \texttt{gemini-embedding-2-preview}.

\noindent\textbf{RAPTOR~\citep{RAPTOR}.}\quad\emph{(i)~Overview.} Recursive abstractive processing: chunks at the leaf level are clustered, each cluster is replaced by an LLM-written summary, and the procedure recurses upward to form a hierarchical tree that lets retrieval pick its abstraction level on the fly. Substrate: cluster tree $\mathcal{T}$ with per-level summaries $\sigma_k$.\quad\emph{(ii)~Per-task implementation.} $\mathcal{T}_1$: $(\hat{\mathcal{E}},\hat{\mathcal{P}},\hat{\mu})\!=\!(\emptyset,\emptyset,\emptyset)$ (the tree never instantiates a symbolic entity layer). $\mathcal{T}_3$: $\hat{a}\!=\!\mathrm{LLM}(q,\mathrm{TreeWalk}_{\phi(q)}(\mathcal{T}))$, descending from the root by cosine relevance until a leaf cluster is reached. $\mathcal{T}_6$: $\mathrm{known}(t)\!=\!\{\sigma_k\!:\!\langle\phi(t),\phi(\sigma_k)\rangle\!\geq\!\theta_{\mathrm{k}}\}$ over the upper levels of $\mathcal{T}$; $\mathrm{gaps}\!=\!\mathrm{sug}\!=\!\emptyset$ since the tree exposes no proactive-emission interface. \emph{Unsupported}: $\mathcal{T}_2$ — no entities to link. $\mathcal{T}_4$ — summaries are text-only paraphrases of text leaves, so $\mu(\sigma_k)\!\equiv\!\{\textit{text}\}$ and $\mathrm{XMEv}\!=\!0$. $\mathcal{T}_5$ — the tree is rebuilt offline; an in-place edit would invalidate every parent summary.\quad\emph{(iii)~Hyperparameters.} chunk size $300$, GMM clustering with branching factor $5$, max tree depth $4$, summary cap $200$ tokens, retrieval top-$k\!=\!10$ at each level, $\theta_{\mathrm{k}}\!=\!0.6$.

\subsection{Graph/Hypergraph RAG}\label{app:b-graph}

Methods here add a text-only entity (or hyper-)graph over the chunk substrate, so $\mathcal{T}_1$ becomes scorable but no cross-modal coreference, lifecycle, or omnimodal evidence is exposed.

\noindent\textbf{GraphRAG~\citep{GraphRAG}.}\quad\emph{(i)~Overview.} Microsoft's graph-based RAG: an LLM extracts entities and relations from each chunk to form a knowledge graph $G\!=\!(V,E)$, the graph is partitioned into Louvain communities $\mathcal{C}_k$ each of which is summarised by a second LLM call, and queries are routed through the relevant community summaries before falling back to entity-level retrieval. Substrate: $G$ together with $\{(\mathcal{C}_k,\sigma_k)\}$.\quad\emph{(ii)~Per-task implementation.} $\mathcal{T}_1$: $\hat{\mathcal{E}}\!=\!V$, $\hat{\mathcal{P}}$ from the extracted relations, $\hat{\mu}\!\equiv\!\{\textit{text}\}$ since extraction runs only on text chunks. $\mathcal{T}_3$: $\hat{a}\!=\!\mathrm{LLM}(q,R)$ with $R\!=\!\bigcup_{k\in\mathrm{TopL}_{\mathcal{C}}}\mathrm{TopK}_{k_v}^{\mathcal{C}_k}\!\langle\phi(q),\phi(v)\rangle$ (community-routed top-$k$). $\mathcal{T}_6$: $\mathrm{known}(t)\!=\!\mathrm{TopL}_{\mathcal{C}_k}\!\langle\phi(t),\phi(\sigma_k)\rangle$ and $(\mathrm{gaps},\mathrm{sug})\!=\!\mathrm{LLM}(t,\{\sigma_k\}_{\mathcal{C}_t})$ by a probe over the selected summaries. \emph{Unsupported}: $\mathcal{T}_2$ — $V$ has no modality tag, so $\hat{\mu}(e_a)\!\cap\!\hat{\mu}(e_b)\!=\!\emptyset$ is vacuously false for every pair. $\mathcal{T}_4$ — without cross-modal entities, $\mu(R)\!\subseteq\!\{\textit{text}\}$ and $\mathrm{XMEv}\!=\!0$. $\mathcal{T}_5$ — $G$ and the Louvain partition are materialised once; an Add/Del/Refresh would invalidate the cached summaries with no recomputation operator.\quad\emph{(iii)~Hyperparameters.} chunk size $300$, Louvain resolution $\gamma\!=\!1.0$, per-community summary $300$ tokens, community top-$L\!=\!5$, entity top-$k_v\!=\!10$, hop expansion off.

\noindent\textbf{LightRAG~\citep{LightRAG}.}\quad\emph{(i)~Overview.} Lightweight dual-level graph retrieval: a low-level index over fine-grained entity-relation pairs is paired with a high-level key-value index over chunk-spanning topical contexts, and queries fetch evidence from both levels in parallel without an explicit community partition. Substrate: dual-level $G\!=\!(V_{\mathrm{lo}}\!\cup\!V_{\mathrm{hi}},\,E)$ with the high level keyed by topic descriptors and valued by chunk anchors.\quad\emph{(ii)~Per-task implementation.} $\mathcal{T}_1$: $\hat{\mathcal{E}}\!=\!V_{\mathrm{lo}}\!\cup\!V_{\mathrm{hi}}$ tagged with the level at which each node was extracted. $\mathcal{T}_3$: $\hat{a}\!=\!\mathrm{LLM}(q,R)$ with $R\!=\!\mathrm{TopK}_{k}^{\mathrm{lo}}\!\langle\phi(q),\phi(v)\rangle\!\cup\!\mathrm{TopK}_{k}^{\mathrm{hi}}\!\langle\phi(q),\phi(v)\rangle$, the union of both rankings. \emph{Unsupported}: $\mathcal{T}_2$ — the graph is text-only; the high level is per-chunk, not per-modality, so no cross-modal pair exists. $\mathcal{T}_4$ — same reason as T2: $\mu(R)\!\subseteq\!\{\textit{text}\}$. $\mathcal{T}_5$ — the dual index is built once and admits no incremental edits. $\mathcal{T}_6$ — there is no community-summary path; the high level is per-chunk rather than per-cluster, so $\mathrm{known}/\mathrm{gaps}$ have no anchor.\quad\emph{(iii)~Hyperparameters.} chunk size $300$, low-level top-$k\!=\!20$, high-level top-$k\!=\!10$, both at cosine threshold $\geq\!0.4$.

\noindent\textbf{HippoRAG2~\citep{HippoRAG2}.}\quad\emph{(i)~Overview.} Neurobiologically-inspired non-parametric continual learning: HippoRAG2 treats the entity graph as a hippocampal-index analogue and propagates query-anchored relevance via Personalized PageRank to mimic pattern completion across associated memories, then continually appends new associations without retraining the encoder. Substrate: $G\!=\!(V,E)$ with per-query PPR vector $\pi(\cdot|q)$ over $V$.\quad\emph{(ii)~Per-task implementation.} $\mathcal{T}_1$: $\hat{\mathcal{E}}\!=\!V$, $\hat{\mu}\!\equiv\!\{\textit{text}\}$. $\mathcal{T}_3$: seeds $\mathcal{Z}_q\!=\!\{v\!\in\!V\!:\!v\,\text{matches a query mention}\}$; $\pi(v|q)\!=\!\alpha\!\cdot\!\mathbf{1}[v\!\in\!\mathcal{Z}_q]\!+\!(1\!-\!\alpha)\!\sum_{u\to v}\pi(u|q)/\deg(u)$ until convergence; $R\!=\!\mathrm{TopK}_{k}\,\pi(\cdot|q)$, $\hat{a}\!=\!\mathrm{LLM}(q,R)$. \emph{Unsupported}: $\mathcal{T}_2$ — single-modality $V$. $\mathcal{T}_4$ — no cross-modal evidence. $\mathcal{T}_5$ — recomputing PPR after every edit is prohibitive, so updates are forbidden. $\mathcal{T}_6$ — no community partition or summaries.\quad\emph{(iii)~Hyperparameters.} chunk size $300$, PPR damping $\alpha\!=\!0.85$, $50$ power iterations, top-$k\!=\!20$, query-entity matcher = exact + lemma.

\noindent\textbf{HyperGraphRAG~\citep{HyperGraphRAG}.}\quad\emph{(i)~Overview.} Hypergraph generalisation of GraphRAG: an LLM extracts $n$-ary relational facts so each hyperedge $h\!\subseteq\!V$ connects all entities of the fact at once, allowing higher-order structure to be retrieved via hyperedge expansion rather than pairwise edges. Substrate: hypergraph $\mathcal{G}_H\!=\!(V,\mathcal{H}^{T})$ with $|h|\!\geq\!2$ for every $h\!\in\!\mathcal{H}^{T}$.\quad\emph{(ii)~Per-task implementation.} $\mathcal{T}_1$: $\hat{\mathcal{E}}\!=\!V$, $\hat{\mathcal{P}}$ enumerates the extracted $n$-ary facts, $\hat{\mu}\!\equiv\!\{\textit{text}\}$. $\mathcal{T}_3$: $R\!=\!\bigcup_{h\in\mathrm{TopK}_{k}\!\langle\phi(q),\phi(h)\rangle}\!\mathrm{Nbr}(h)$ where $\mathrm{Nbr}(h)\!=\!\bigcup_{v\in h}\{v\}$; $\hat{a}\!=\!\mathrm{LLM}(q,R)$. \emph{Unsupported}: $\mathcal{T}_2$ — the hypergraph is text-only; $\hat{\mu}$ collapses to a single value. $\mathcal{T}_4$ — same as T2: $\mathrm{XMEv}\!=\!0$. $\mathcal{T}_5$ — no lifecycle agent; deletion would leave dangling hyperedges. $\mathcal{T}_6$ — no community layer over the hypergraph.\quad\emph{(iii)~Hyperparameters.} chunk size $300$, max hyperedge arity $n\!=\!5$, top-$k\!=\!20$ over hyperedges, expansion depth $1$ (hyperedge $\to$ vertices), no community summaries.

\subsection{KG Construction}\label{app:b-kgconst}

Schema-grounded KG builders that constrain extraction to a typed ontology, supporting precise $\mathcal{T}_3$ but no cross-modal or routing layer.

\noindent\textbf{iText2KG~\citep{iText2KG}.}\quad\emph{(i)~Overview.} Iterative text-to-knowledge-graph construction: an LLM extracts entity-relation triples while a developing ontology $\mathcal{O}$ is refined as new documents arrive, producing a high-precision typed graph at the cost of lifecycle flexibility. Substrate: ontology-grounded $G\!=\!(V,E)$ pinned to schema $\mathcal{O}$.\quad\emph{(ii)~Per-task implementation.} $\mathcal{T}_1$: $\hat{\mathcal{E}}\!=\!V$ typed by $\mathcal{O}$, $\hat{\mathcal{P}}$ enumerates ontology-typed triples $(s,r,o)$, $\hat{\mu}\!\equiv\!\{\textit{text}\}$. $\mathcal{T}_3$: $\hat{a}\!=\!\mathrm{LLM}(q,\mathrm{Solve}(\mathrm{SQ}(q),G))$ where $\mathrm{SQ}$ first rewrites $q$ into a typed graph query against $\mathcal{O}$. $\mathcal{T}_5$ \emph{partial}: only $\textsc{Add}(f)$ is operationalised --- it triggers an ontology rebuild $\mathcal{O}\!\to\!\mathcal{O}'$ and re-extraction over the new asset, while Del/Refresh require manual intervention; T5 is therefore evaluated only on Add-only sub-streams of the lifecycle probe set. \emph{Unsupported}: $\mathcal{T}_2$ — there is no cross-modal coreference test $\chi$; the ontology has no modality slot. $\mathcal{T}_4$ — without $\chi$ and with text-only extraction, $\mathrm{XMEv}\!=\!0$. $\mathcal{T}_6$ — the ontology defines types but no community summaries.\quad\emph{(iii)~Hyperparameters.} chunk size $300$, ontology max depth $3$, triple confidence threshold $0.7$, $\mathrm{SQ}$ compiled to Cypher, query timeout $5$\,s.

\noindent\textbf{OntoKG~\citep{OntoKG}.}\quad\emph{(i)~Overview.} Ontology-grounded KG construction with a frozen domain schema: extraction is constrained to ontology-validated triples for high precision, but the schema cannot evolve and the graph cannot be edited after build. Substrate: fixed-ontology $G$ pinned to a hand-curated $\mathcal{O}$.\quad\emph{(ii)~Per-task implementation.} $\mathcal{T}_1$: $\hat{\mathcal{E}}\!=\!V$ typed by $\mathcal{O}$, $\hat{\mathcal{P}}$ as schema-validated triples; entries that fail ontology validation are dropped before scoring. $\mathcal{T}_3$: $\hat{a}\!=\!\mathrm{LLM}(q,\mathrm{SPARQL}(q,G))$, where the LLM compiles $q$ to a structured SPARQL-like query against $\mathcal{O}$. \emph{Unsupported}: $\mathcal{T}_2$ — the ontology has no modality slot, $\hat{\mu}$ is undefined. $\mathcal{T}_4$ — text-only triples, $\mathrm{XMEv}\!=\!0$. $\mathcal{T}_5$ — the schema is frozen, so any $\Delta_n$ is rejected at validation. $\mathcal{T}_6$ — the schema defines types but not community summaries.\quad\emph{(iii)~Hyperparameters.} chunk size $300$, schema-confidence threshold $\geq\!0.6$, SPARQL query timeout $5$\,s, three-shot SPARQL compiler prompts.

\subsection{Agentic KG}\label{app:b-akg}

Live KG engines with partial CRUD and operator-driven retrieval, but no cross-modal or multi-angle layer.

\noindent\textbf{KAG/OpenSPG~\citep{KAG}.}\quad\emph{(i)~Overview.} Production-grade open-source KG engine: combines schema-driven extraction with operator-driven retrieval and partial CRUD on a live typed KG, designed as the back-end of OpenSPG-powered domain assistants. Substrate: live typed KG $G$ with first-class Add and Del operators.\quad\emph{(ii)~Per-task implementation.} $\mathcal{T}_1$: $\hat{\mathcal{E}}\!=\!V$ with full type tags from the OpenSPG schema, $\hat{\mathcal{P}}$ as typed triples, $\hat{\mu}\!\equiv\!\{\textit{text}\}$. $\mathcal{T}_3$: anchor $v_q\!=\!\arg\max_{v\in V}\!\langle\phi(q),\phi(v)\rangle$, $R\!=\!\mathrm{Walk}_{\ell}(v_q;G)$ via typed edge expansion, $\hat{a}\!=\!\mathrm{LLM}(q,R)$. $\mathcal{T}_5$ \emph{partial}: Add and Del are first-class operators with referential-integrity checks; Refresh is implemented as Del+Add and therefore loses $\phi(h_{\mathrm{old}})$'s neighbourhood embedding (no atomic in-place rewrite), reducing the refresh-F1 component of Eq.~\eqref{eq:m-t5}. \emph{Unsupported}: $\mathcal{T}_2$ — KG nodes are text-only; the schema lacks a modality slot. $\mathcal{T}_4$ — no cross-modal evidence. $\mathcal{T}_6$ — no community-summary materialisation.\quad\emph{(iii)~Hyperparameters.} chunk size $300$, walk depth $\ell\!=\!2$, KG capacity $50$\,k entities, Refresh as Del+Add, anchor matching = lemma + cosine top-$1$.

\noindent\textbf{Graphiti~\citep{Graphiti}.}\quad\emph{(i)~Overview.} Bi-temporal KG for autonomous agents: every edge carries a validity interval, so the graph supports time-aware queries and version-controlled writes that preserve historical state instead of overwriting it. Substrate: temporal graph $G_t\!=\!(V,E_t)$ with per-edge timestamps.\quad\emph{(ii)~Per-task implementation.} $\mathcal{T}_1$: $\hat{\mathcal{E}}\!=\!V$, $\hat{\mathcal{P}}$ enumerates temporally-stamped triples; $\hat{\mu}\!\equiv\!\{\textit{text}\}$. $\mathcal{T}_3$: $R\!=\!\mathrm{TWalk}_{\ell}(v_q;t_q)\!=\!\{(u,v,t)\!\in\!E_t\!:\,t\!\leq\!t_q\,\wedge\,d_G(v_q,v)\!\leq\!\ell\}$ where $t_q$ is the query reference time. $\mathcal{T}_5$ \emph{partial}: Add stamps a new edge with current time; Del marks the old as expired but does \emph{not} cascade-evict orphan entities (so $\mathcal{S}(e)\!=\!\emptyset$ entities can persist); Refresh writes a new versioned edge while keeping $h_{\mathrm{old}}$ retrievable, which prevents the refresh-F1 from cleanly capturing the new fact. $\mathcal{T}_6$: $\mathrm{known}(t_q)\!=\!\{\sigma_k(t_q)\}$ over time-conditioned clusters, $\mathrm{gaps}\!=\!\mathrm{sug}\!=\!\emptyset$ (no proactive emission). \emph{Unsupported}: $\mathcal{T}_2$ — text-only nodes, no $\chi$. $\mathcal{T}_4$ — text-only evidence, $\mathrm{XMEv}\!=\!0$.\quad\emph{(iii)~Hyperparameters.} chunk size $300$, temporal walk depth $\ell\!=\!2$, retention horizon $=$ full history, $t_q$ = corpus snapshot time of the query, cluster summary cap $200$ tokens.

\subsection{Memory Agents}\label{app:b-mem}

Memory-oriented agents whose substrate is a flat or structured memory store with full or partial CRUD, but no symbolic graph layer over $\mathcal{M}$.

\noindent\textbf{MemGPT/Letta~\citep{MemGPT}.}\quad\emph{(i)~Overview.} Hierarchical-memory agent: an LLM manages a tiered memory system (main context, recall buffer, archival store) via function calls, paging information in and out of the context window so a fixed-size LLM can reason over arbitrarily long histories. Substrate: hierarchical memory $\{m_i\}$ partitioned into main, recall, and archival layers.\quad\emph{(ii)~Per-task implementation.} $\mathcal{T}_1$: $\hat{\mathcal{E}}$ approximates memory entries (no NER pass, the memory store contains free-form blobs written by the agent), $\hat{\mathcal{P}}\!=\!\emptyset$, $\hat{\mu}\!\equiv\!\{\textit{text}\}$. $\mathcal{T}_3$: $R\!=\!\mathrm{ANN}_{k}(\phi(q),\{\phi(m_i)\})$ over the recall + archival layers, with the main layer always concatenated. $\mathcal{T}_5$: tool-call CRUD via the LLM's function-call interface --- $\textsc{Add},\textsc{Del},\textsc{Refresh}$ each map to a single tool invocation, but no cross-memory cascade or orphan detection is run, so a Del that breaks references is silently committed. \emph{Unsupported}: $\mathcal{T}_2$ — no entity graph and no modality typing on memory blobs. $\mathcal{T}_4$ — same: $\mu(m_i)\!\equiv\!\{\textit{text}\}$. $\mathcal{T}_6$ — no clustering or community layer over the memory store.\quad\emph{(iii)~Hyperparameters.} main memory $8$\,k tokens, recall memory $100$\,k tokens, archival store unbounded, ANN top-$k\!=\!10$, tool-call budget $20$ per query.

\noindent\textbf{Mem0~\citep{Mem0}.}\quad\emph{(i)~Overview.} Production-grade memory layer for LLM agents: stores user interactions in a hybrid temporal-semantic memory tree with first-class CRUD APIs, optimised for personalised assistant deployments where recency and relevance both matter. Substrate: structured memory tree $\mathcal{T}_m$ with semantic and temporal indexes.\quad\emph{(ii)~Per-task implementation.} $\mathcal{T}_1$: memory nodes act as quasi-entities; $\hat{\mathcal{E}}\!=\!\{m_i\}$, $\hat{\mathcal{P}}$ from the LLM-generated ``fact'' field of each node, $\hat{\mu}\!\equiv\!\{\textit{text}\}$. $\mathcal{T}_3$: $R\!=\!\mathrm{TopK}_{k}\bigl(\alpha\!\cdot\!\langle\phi(q),\phi(m)\rangle\!+\!(1\!-\!\alpha)\!\cdot\!\mathrm{Recency}(m)\bigr)$ blends semantic similarity and temporal decay. $\mathcal{T}_5$: full CRUD on $\mathcal{T}_m$ --- Add appends a new memory, Del removes the targeted node, and Update rewrites the ``fact'' field in place; the tree's parent-child semantics handle simple cascades but offer no $\chi$-aware coreference. \emph{Unsupported}: $\mathcal{T}_2$ — no graph layer, no modality typing. $\mathcal{T}_4$ — same: text-only memory. $\mathcal{T}_6$ — the tree exposes no community partition or summary layer.\quad\emph{(iii)~Hyperparameters.} hybrid weight $\alpha\!=\!0.7$, top-$k\!=\!15$, recency half-life $7$ days, max tree depth $5$.

\noindent\textbf{A-MEM~\citep{A-MEM}.}\quad\emph{(i)~Overview.} Agentic memory with structured slots: each interaction the agent writes a tuple of (entity, relation, evidence) slots into typed memory cells, supporting both lookup and a limited gap-detection probe driven by slot occupancy statistics. Substrate: agentic memory $m$ with $\{(e_i,r_i,\text{evidence}_i)\}$ slot tuples.\quad\emph{(ii)~Per-task implementation.} $\mathcal{T}_1$: structured memory entries decoded as $\hat{\mathcal{E}},\hat{\mathcal{P}}$ from the entity and relation slots; $\hat{\mu}\!\equiv\!\{\textit{text}\}$. $\mathcal{T}_3$: $R\!=\!\mathrm{Slot\text{-}Match}(q,\{m_i\})$ where $\mathrm{Slot\text{-}Match}$ retrieves entries whose entity slot matches a query mention and whose relation slot is semantically close to the query predicate. $\mathcal{T}_5$: full CRUD on slots --- Add appends a slot tuple, Del removes one, Update rewrites individual fields; supports per-slot consistency checks but not graph-wide cascade. $\mathcal{T}_6$ \emph{partial}: an LLM emits gaps from slot-occupancy statistics ($r_i$ rare in the corpus $\Rightarrow$ candidate gap), but $\mathrm{known}/\mathrm{sug}$ are not differentiated --- the probe collapses both into the same returned list, halving its T6 score relative to a community-grounded probe. \emph{Unsupported}: $\mathcal{T}_2$ — slot tuples have no modality field, so no cross-modal pair exists. $\mathcal{T}_4$ — same: $\mu(R)\!\subseteq\!\{\textit{text}\}$.\quad\emph{(iii)~Hyperparameters.} $8$ slots per entry, retrieval threshold $\geq\!0.5$, gap probe top-$k\!=\!5$, occupancy bin width $1{\%}$.

\subsection{Omnimodal RAG}\label{app:b-omni}

Index-per-modality methods that fuse retrieval across $\mathcal{M}$ but drop the lifecycle and discovery layers.

\noindent\textbf{OmniRAG-Agent~\citep{OmniRAG-Agent}.}\quad\emph{(i)~Overview.} Per-modality RAG with hashing-based cross-modal linking: each modality gets its own cosine index, queries fan out to all modalities in parallel, and modality fusion happens at the prompt level by concatenating the per-modality top-$k$. Substrate: per-modality indexes $\{I_m\}_{m\in\mathcal{M}}$ with a shared LSH-based linker over $\phi$.\quad\emph{(ii)~Per-task implementation.} $\mathcal{T}_1$: $\hat{\mathcal{E}}\!=\!\bigcup_m\hat{\mathcal{E}}_m$ from per-modality NER; $\hat{\mu}(e)\!=\!\{m\!:\!e\!\in\!\hat{\mathcal{E}}_m\}$ records the observed modalities. $\mathcal{T}_2$: $\hat{L}\!=\!\{(e_a,e_b)\!:\!\mathrm{LSH}_{\theta}(\phi(e_a),\phi(e_b))\!=\!1\}$, hashing-based linking with no LLM verifier --- semantically aligned cross-modal mentions pass when their hash bands collide, but unrelated entities can also pass on accidental collisions, so T2 receives a higher false-positive rate than $\chi$-tested methods. $\mathcal{T}_3$: $R\!=\!\bigcup_m\mathrm{TopK}_{k_m}^{I_m}\langle\phi(q),\phi(c)\rangle$, per-modality top-$k$ then merge. $\mathcal{T}_4$: $\hat{a}\!=\!\mathrm{LLM}(q,\bigoplus_m R_m)$, modality fusion by ordered concatenation; $\mu(R)\!=\!\bigcup_m\{m\}\!=\!\mathcal{M}$ so $\mathrm{XMEv}$ is at ceiling whenever fusion happens. \emph{Unsupported}: $\mathcal{T}_5$ — the per-modality indexes admit no incremental $\Delta_n$; updates require offline rebuild. $\mathcal{T}_6$ — no community partition over the multi-index substrate.\quad\emph{(iii)~Hyperparameters.} chunk size $300$ for text, CLIP for image embedding, Whisper-v3 for audio transcription, per-modality top-$k_m\!=\!8$, LSH bands $b\!=\!8$, rows per band $r\!=\!4$, threshold $\theta\!=\!0.7$.

\subsection{Agentic KE}\label{app:b-ake}

End-to-end agentic systems that combine an entity graph with memory and partial omnimodal coverage; they support every task at least partially but with ad-hoc components rather than the orthogonal C1/C2/C3 design.

\noindent\textbf{OpenClaw~\citep{OpenClaw}.}\quad\emph{(i)~Overview.} Anchor system for agentic knowledge engineering: combines an omnimodal entity graph with a flat memory store inside a closed-loop perception-reasoning-tool-use agent, unifying perception, reasoning, tool use, and memory in one pipeline over personal file collections. We use OpenClaw as the $\Delta$ baseline of Table~\ref{tab:main}. Substrate: graph-plus-memory $\mathcal{G}^{\circ}\!=\!(V,E,M)$ with omnimodal extraction.\quad\emph{(ii)~Per-task implementation.} $\mathcal{T}_1$: $\hat{\mathcal{E}}\!=\!V$ with $\hat{\mu}(e)$ filled by a per-modality entity-extraction agent (one extractor invoked per modality of the source asset). $\mathcal{T}_2$: $\hat{L}\!=\!\{(e_a,e_b)\!:\!\mathrm{Linker}_{\mathrm{LLM}}(e_a,e_b)\!=\!1\}$, an LLM-style linking pass without the cosine pre-filter of Eq.~\eqref{eq:coref}, which raises both false positives (unverified semantic agreement) and runtime cost. $\mathcal{T}_3$: $R\!=\!\mathrm{TopK}_{k}\!\langle\phi(q),\phi(v)\rangle\!\cup\!M_q$, top-$k$ over $V$ unioned with the relevant memory slice. $\mathcal{T}_4$: $\hat{a}\!=\!\mathrm{LLM}(q,\bigoplus_m R_m)$, modality concat over the modalities present in $R$. $\mathcal{T}_5$ \emph{partial}: Add and Del are committed by the agent; Refresh only re-extracts the affected file rather than rewriting $h_{\mathrm{old}}\!\to\!h_{\mathrm{new}}$ in place, so $\phi(h_{\mathrm{new}})$ drifts and downstream retrieval lags. $\mathcal{T}_6$: $(\mathrm{known},\mathrm{gaps},\mathrm{sug})\!=\!\mathrm{LLM}(t,V)$ probe without community routing --- a single prompt against the global graph emits all three lists, which under-emphasises $\mathrm{gaps}$ versus $\mathrm{known}$ when $V$ is large.\quad\emph{(iii)~Hyperparameters.} chunk size $300$, top-$k\!=\!15$, memory cap $10$\,k entries, refresh policy $=$ re-extract affected file, linker prompt = $3$-shot, OpenClaw-base LLM agent.

\noindent\textbf{Cognee~\citep{Cognee}.}\quad\emph{(i)~Overview.} Schema-aware variant of OpenClaw: layers a Cognee ontology over the omnimodal graph so every extracted entity carries a typed label and edges carry typed relations, trading some recall for higher precision on typed retrieval. Substrate: typed-relation $\mathcal{G}^{\circ,\tau}\!=\!(V,E,M)$ with type assignment $\tau\!:\!V\!\to\!\mathcal{O}_{\mathrm{Cognee}}$.\quad\emph{(ii)~Per-task implementation.} $\mathcal{T}_1$: $\hat{\mathcal{E}}\!=\!V$ tagged with Cognee ontology labels; entries that fail ontology validation are dropped. $\mathcal{T}_3$: $R\!=\!\mathrm{TopK}_{k}\!\bigl(\,\alpha\!\cdot\!\langle\phi(q),\phi(v)\rangle\!+\!(1\!-\!\alpha)\!\cdot\!\mathbf{1}[\tau(v)\!=\!\tau^{*}(q)]\,\bigr)$ where $\tau^{*}(q)$ is the inferred query type, biasing retrieval toward same-type neighbours. $\mathcal{T}_2$, $\mathcal{T}_4$, $\mathcal{T}_5$, $\mathcal{T}_6$ inherit OpenClaw's mechanisms verbatim.\quad\emph{(iii)~Hyperparameters.} chunk size $300$, type-bias $\alpha\!=\!0.6$, top-$k\!=\!15$, ontology depth $2$, all other settings as OpenClaw.

\noindent\textbf{LLM Wiki~\citep{LLMWiki}.}\quad\emph{(i)~Overview.} Wiki-style materialisation of OpenClaw: each entity is rendered as a canonical multi-section page (text, image, audio, video, related), which becomes the retrieval unit and gives the LLM a Wikipedia-like view of the corpus. Substrate: page store $\mathcal{P}_{\mathrm{wiki}}\!=\!\{p_e\!:\!e\!\in\!V\}$ over the same omnimodal $\mathcal{G}^{\circ}$.\quad\emph{(ii)~Per-task implementation.} $\mathcal{T}_1$: identical to OpenClaw. $\mathcal{T}_3$: $R\!=\!\mathrm{TopK}_{k}\!\langle\phi(q),\phi(p_e)\rangle$ where $\phi(p_e)$ embeds the page rather than the bare entity, so retrieval is page-level and the LLM sees structured sections instead of raw nodes. $\mathcal{T}_4$: $\hat{a}\!=\!\mathrm{LLM}(q,\bigoplus_e \bigoplus_m \mathrm{Sec}_m(p_e))$, concatenating each retrieved page's per-modality sections in fixed order. $\mathcal{T}_6$: probe operates over page summaries instead of the raw graph, $(\mathrm{known},\mathrm{gaps},\mathrm{sug})\!=\!\mathrm{LLM}(t,\{\mathrm{summary}(p_e)\}_{e\in\mathrm{TopL}})$. $\mathcal{T}_2$ and $\mathcal{T}_5$ inherit OpenClaw verbatim.\quad\emph{(iii)~Hyperparameters.} page max $2$\,k tokens, per-modality section cap $400$ tokens, retrieval top-$k\!=\!8$ pages, summary cap $200$ tokens, all other settings as OpenClaw.

\noindent\textbf{Graphify~\citep{Graphify}.}\quad\emph{(i)~Overview.} Property-graph-first variant of OpenClaw: nodes carry typed properties (e.g., \texttt{birth\_year}, \texttt{file\_format}) and edges carry typed relations, so retrieval can filter on properties before semantic top-$k$, mirroring Neo4j-style index-and-then-rank query plans. Substrate: property graph $\mathcal{G}^{\circ,\mathrm{prop}}$ on top of the omnimodal $\mathcal{G}^{\circ}$ of OpenClaw.\quad\emph{(ii)~Per-task implementation.} $\mathcal{T}_1$: as OpenClaw, with each $e\!\in\!V$ additionally annotated by a sparse property dictionary $\mathrm{prop}(e)$. $\mathcal{T}_3$: $V_q\!=\!\sigma_{\mathrm{prop}(q)}(V)\!=\!\{v\!:\!\mathrm{prop}(v)\,\text{satisfies the parsed query constraints}\}$; $R\!=\!\mathrm{TopK}_{k}\!\langle\phi(q),\phi(v)\rangle$ over $V_q$ rather than the full $V$, pruning off-topic types before semantic ranking. $\mathcal{T}_2$, $\mathcal{T}_4$, $\mathcal{T}_5$, $\mathcal{T}_6$ otherwise match OpenClaw.\quad\emph{(iii)~Hyperparameters.} chunk size $300$, property index over the top-$5$ most frequent attributes, post-filter top-$k\!=\!15$, all other settings as OpenClaw.

\subsection{KnowledgeClaw Ablations}\label{app:b-abl}

The three internal ablations remove exactly one contribution to isolate its empirical effect; notation is identical to KnowledgeClaw (\S\ref{app:b-kc}).

\noindent\textbf{KC w/o C1.}\quad\emph{(i)~Overview.} KnowledgeClaw without the omnimodal hypergraph: the substrate degenerates to a text-only entity graph extracted by $\mathrm{Extract}_{\mathrm{LLM}}$ over text chunks alone, with neither within-modality canonicalisation $\kappa$ nor cross-modal coreference $\chi$. C2 and C3 still run, but on this strictly weaker substrate. Substrate: text-only $\mathcal{G}_{\textit{text}}\!\subset\!\mathcal{G}$.\quad\emph{(ii)~Per-task implementation.} $\mathcal{T}_1$: $\hat{\mathcal{E}}_{\textit{text}}\!\subseteq\!\mathcal{E}$ extracted from text mentions only, $\hat{\mu}\!\equiv\!\{\textit{text}\}$, so the modality-coverage component $C_\mu$ of Eq.~\eqref{eq:m-t1} is fixed at $1/|\mathcal{M}|$. $\mathcal{T}_3$: RRF + $\ell$-hop on the text-restricted graph exactly as in Eq.~\eqref{eq:sub}, but with $|\mathcal{E}\!\cup\!\mathcal{H}|$ shrunk to its text-only subset, so multi-hop chains that previously passed through image, audio, or video bridges are unreachable. $\mathcal{T}_5$: Build--Query--Reason--Update with cascade-delete (Eq.~\eqref{eq:cascade}) restricted to text orphans; refresh works in place since $\phi$ is unchanged. $\mathcal{T}_6$: $\mathrm{LLM}_{\mathrm{disc}}$ on text-only communities (Eq.~\eqref{eq:disc-route}), with the gap-emission constrained to the text vocabulary. \emph{Unsupported}: $\mathcal{T}_2$ — without $\chi$ the cross-modal coreference test cannot fire, $\hat{L}\!=\!\emptyset$. $\mathcal{T}_4$ — every $\mu^{*}(q)$ that includes a non-text modality forces $\mathrm{XMEv}(R)\!=\!0$ since $\mu(R)\!\subseteq\!\{\textit{text}\}$. The empirical drop of $-30.4$ OKE in Table~\ref{tab:main} confirms C1 as the structural foundation that the other two contributions build on.\quad\emph{(iii)~Hyperparameters.} same as KnowledgeClaw, except extraction is restricted to the \texttt{text} field of $\rho(\mathcal{D})$, $\theta_{\chi}$ is undefined (no $\chi$ pass), and the bipartite incidence matrix is restricted to $\hat{\mathcal{E}}_{\textit{text}}$.

\noindent\textbf{KC w/o C2.}\quad\emph{(i)~Overview.} KnowledgeClaw without the closed-loop lifecycle agent: C1 builds the full omnimodal hypergraph and C3 serves it, but every corpus event $\Delta_n$ is rejected at runtime so the graph is frozen after the initial build. Substrate: full $\mathcal{G}$ from Algorithm~\ref{alg:kc-c1}, but no Build--Query--Reason--Update operator.\quad\emph{(ii)~Per-task implementation.} $\mathcal{T}_1$: full $\mathrm{Extract}\!+\!\kappa$ as in Algorithm~\ref{alg:kc-c1}. $\mathcal{T}_2$: $\chi$ test (Eq.~\eqref{eq:coref}) with cosine pre-filter and LLM verifier. $\mathcal{T}_3$: RRF + $\ell$-hop (Eq.~\eqref{eq:sub}). $\mathcal{T}_4$: $\chi$-aware subgraph retrieval --- the seed set is restricted to nodes with $\hat{\mu}(R)\!\supseteq\!\mu^{*}(q)$ before $\ell$-hop expansion. $\mathcal{T}_6$: $\mathrm{LLM}_{\mathrm{disc}}$ over communities (Eq.~\eqref{eq:disc-route}). \emph{Unsupported}: $\mathcal{T}_5$ — Algorithm~\ref{alg:kc-c2} is disabled, so $\Delta_n$ events queued by the lifecycle probe set are rejected before \textsc{Build} ever runs; the NL-CRUD verification therefore evaluates the post-event response against a stale KB (the keyword check fails on every \textsc{Del}/\textsc{Refresh}), and fact-refresh keeps returning $h_{\mathrm{old}}$, so both halves of Eq.~\eqref{eq:m-t5} score zero. T1--T4 and T6 are unaffected, isolating the lifecycle effect at $-7.8$ OKE.\quad\emph{(iii)~Hyperparameters.} identical to KnowledgeClaw; the $\Delta_n$ queue is logged for inspection but never consumed by a \textsc{Build} pass.

\noindent\textbf{KC w/o C3.}\quad\emph{(i)~Overview.} KnowledgeClaw without multi-angle routing: C1 and C2 still run, but the application layer collapses to a single flat-cosine retrieval pass --- no RRF, no $\ell$-hop expansion, no community summary, no $\mathrm{LLM}_{\mathrm{disc}}$. Substrate: full $\mathcal{G}$ + lifecycle from Algorithms~\ref{alg:kc-c1}--\ref{alg:kc-c2}, but Algorithm~\ref{alg:kc-c3} is replaced by flat top-$k$.\quad\emph{(ii)~Per-task implementation.} $\mathcal{T}_1$: full $\mathrm{Extract}\!+\!\kappa$. $\mathcal{T}_2$: $\chi$ test. $\mathcal{T}_3$ and $\mathcal{T}_4$: $R\!=\!\mathrm{TopK}_{k}\!\langle\phi(q),\phi(\mathcal{E}\!\cup\!\mathcal{H})\rangle$ --- a single-pass cosine without RRF, without $\ell$-hop expansion, and without modality routing, so multi-hop chains cannot be assembled and $\mathrm{XMEv}$ is no longer guaranteed; T3 drops to $24.0$ and T4 to $35.5$ in Table~\ref{tab:main}. $\mathcal{T}_5$: B--Q--R--U as in Algorithm~\ref{alg:kc-c2}, fully intact, so T5 is preserved. $\mathcal{T}_6$: $\mathrm{known}(t)\!=\!\mathrm{TopK}_{k}\!\langle\phi(t),\phi(\mathcal{E})\rangle$ over the bare entity index, with $\mathrm{gaps}\!=\!\mathrm{sug}\!=\!\emptyset$ (no proactive emission since $\mathrm{LLM}_{\mathrm{disc}}$ is removed); T6 drops to $48.1$. \emph{Unsupported}: none architecturally, but T3/T4/T6 collapse to flat retrieval, isolating the application-layer effect at $-7.9$ OKE.\quad\emph{(iii)~Hyperparameters.} flat top-$k\!=\!20$ for $\mathcal{T}_3,\mathcal{T}_4,\mathcal{T}_6$; the community partition is still computed at build time but never consulted at query time; all other settings as KnowledgeClaw.

\subsection{KnowledgeClaw}\label{app:b-kc}

KnowledgeClaw composes C1 (Algorithm~\ref{alg:kc-c1}), C2 (Algorithm~\ref{alg:kc-c2}), and C3 (Algorithm~\ref{alg:kc-c3}) into the unique system that scores on every $\mathcal{T}_j$ and therefore carries no ``-'' in Table~\ref{tab:design}.

\noindent\textbf{KnowledgeClaw.}\quad\emph{(i)~Overview.} The proposed method: C1 (Algorithm~\ref{alg:kc-c1}) materialises a typed omnimodal hypergraph $\mathcal{G}$, C2 (Algorithm~\ref{alg:kc-c2}) maintains it under arbitrary $\Delta_n$ event streams, and C3 (Algorithm~\ref{alg:kc-c3}) serves it through three calibrated retrieval angles routed by a lightweight LLM gate. The three contributions compose as $\mathrm{C1}\!\to\!\mathrm{C2}^{*}\!\to\!\mathrm{C3}$: a one-time C1 pass yields $\mathcal{G}_0$, C2 runs once per event $\Delta_n$ producing $(\mathcal{G}_{n+1},\mathcal{L}_{n+1})$, and serving issues C3 against the latest $\mathcal{G}_n$ without blocking on C2 since the bipartite incidence matrix is copy-on-write. Substrate: $\mathcal{G}\!=\!(\mathcal{E},\mathcal{H},\mathcal{A},\mathcal{P};\mu,\tau,\chi,\phi)$ from Eq.~\eqref{eq:hg}.\quad\emph{(ii)~Per-task implementation.} $\mathcal{T}_1$: $(\hat{\mathcal{E}},\hat{\mathcal{P}},\hat{\mu})\!=\!\mathrm{Extract}_{\mathrm{LLM}}$ over each $f\!\in\!\mathcal{A}$ followed by within-modality canonicalisation $\kappa$. $\mathcal{T}_2$: cross-modal $\chi$ test of Eq.~\eqref{eq:coref} with cosine pre-filter at $\theta_{\chi}\!=\!0.78$ and an LLM verifier over $\mathcal{S}(e_a)\!\cup\!\mathcal{S}(e_b)$. $\mathcal{T}_3$: subgraph retrieval $R_{\mathrm{sub}}$ via reciprocal-rank fusion plus $\ell$-hop walk (Eq.~\eqref{eq:sub}). $\mathcal{T}_4$: same $R_{\mathrm{sub}}$ optionally concatenated with $R_{\mathrm{com}}$ when the router fires both, $\hat{a}\!=\!\mathrm{LLM}_{\mathrm{gen}}(q,\bigoplus_{a\in\pi^{*}(q)}R_a)$ (Eq.~\eqref{eq:disc-route}); the cited modalities $\hat{\mu}(R)$ cover $\mu^{*}(q)$ by construction since each $h\!\in\!\mathcal{H}$ records its supporting assets $\mathcal{A}_h$, so $\mathrm{XMEv}(R)\!=\!1$ on every gold-cross-modal query. $\mathcal{T}_5$: Build--Query--Reason--Update (Eq.~\eqref{eq:loop}) per Algorithm~\ref{alg:kc-c2}, with cascade delete (Eq.~\eqref{eq:cascade}) and in-place refresh (Eq.~\eqref{eq:refresh}); the audit log $\mathcal{L}_n$ replays every event byte-for-byte. $\mathcal{T}_6$: $\mathrm{LLM}_{\mathrm{disc}}$ over the topic-conditioned community summaries $\{\sigma_k\}_{\mathcal{C}_t^{\mathrm{com}}}$ (Eq.~\eqref{eq:disc-route}), emitting $\mathrm{known}/\mathrm{gaps}/\mathrm{sug}$ in a single call. \emph{Unsupported:} none --- every $\mathcal{T}_j$ has a closed-form implementation grounded in C1--C3, which is what makes KnowledgeClaw the only entry whose row carries no ``-'' in Table~\ref{tab:design}.\quad\emph{(iii)~Hyperparameters.} chunk size $300$, RRF top-$K_{\mathcal{E}}\!=\!K_{\mathcal{H}}\!=\!40$, walk depth $\ell\!=\!2$, RRF constant $k_0\!=\!60$, coreference threshold $\theta_{\chi}\!=\!0.78$, Louvain resolution $\gamma\!=\!1.0$, router gate $\theta_{\pi}\!=\!0.5$, discovery threshold $\theta_{\mathrm{k}}\!=\!0.6$, all temperatures in $[0,0.4]$ on \texttt{gemini-3.1-flash-lite-preview} with \texttt{gemini-embedding-2-preview} for $\phi$; the community partition is recomputed once per build while RRF and router caches are invalidated only on a committed $\Delta_n$. The composition is both ablation- and pipeline-orthogonal --- each contribution can be cached or scaled independently --- yielding the \S\ref{5.7} profile of \$1.95/persona, $41$\,min build, and $580$\,ms serving while ranking first on every $\mathcal{T}_j$ at OKE $59.8$ ($+13.6$ over LLM Wiki), showing that the C1+C2+C3 composition is necessary and sufficient for jointly satisfying properties (i)--(iii).

\clearpage


%%%%%%%%%%%%%%%%%%%%%%%%%%%%%%%%%%%%%%%%%%%%%%%%%%%%%%%%%%%%

\newpage
\section*{NeurIPS Paper Checklist}

\begin{enumerate}

\item {\bf Claims}
    \item[] Question: Do the main claims made in the abstract and introduction accurately reflect the paper's contributions and scope?
    \item[] Answer: \answerYes{}
    \item[] Justification: The abstract and Section~\ref{sec:intro} state that KnowledgeClaw is the first agentic KE system jointly satisfying omnimodality, closed-loop lifecycle, and multi-angle application, and Section~\ref{sec:exp} with Table~\ref{tab:main} verify these claims by ranking first on all six OmniKE tasks at OKE 59.8 (+13.6 over the strongest baseline).
    \item[] Guidelines:
    \begin{itemize}
        \item The answer \answerNA{} means that the abstract and introduction do not include the claims made in the paper.
        \item The abstract and/or introduction should clearly state the claims made, including the contributions made in the paper and important assumptions and limitations. A \answerNo{} or \answerNA{} answer to this question will not be perceived well by the reviewers.
        \item The claims made should match theoretical and experimental results, and reflect how much the results can be expected to generalize to other settings.
        \item It is fine to include aspirational goals as motivation as long as it is clear that these goals are not attained by the paper.
    \end{itemize}

\item {\bf Limitations}
    \item[] Question: Does the paper discuss the limitations of the work performed by the authors?
    \item[] Answer: \answerYes{}
    \item[] Justification: Appendix~F (Limitations and Future Work) explicitly discusses scope of persona/modality coverage, dependence on external LLM and embedding APIs, and assumptions underlying the lifecycle and routing claims.
    \item[] Guidelines:
    \begin{itemize}
        \item The answer \answerNA{} means that the paper has no limitation while the answer \answerNo{} means that the paper has limitations, but those are not discussed in the paper.
        \item The authors are encouraged to create a separate ``Limitations'' section in their paper.
        \item The paper should point out any strong assumptions and how robust the results are to violations of these assumptions (e.g., independence assumptions, noiseless settings, model well-specification, asymptotic approximations only holding locally). The authors should reflect on how these assumptions might be violated in practice and what the implications would be.
        \item The authors should reflect on the scope of the claims made, e.g., if the approach was only tested on a few datasets or with a few runs. In general, empirical results often depend on implicit assumptions, which should be articulated.
        \item The authors should reflect on the factors that influence the performance of the approach. For example, a facial recognition algorithm may perform poorly when image resolution is low or images are taken in low lighting. Or a speech-to-text system might not be used reliably to provide closed captions for online lectures because it fails to handle technical jargon.
        \item The authors should discuss the computational efficiency of the proposed algorithms and how they scale with dataset size.
        \item If applicable, the authors should discuss possible limitations of their approach to address problems of privacy and fairness.
        \item While the authors might fear that complete honesty about limitations might be used by reviewers as grounds for rejection, a worse outcome might be that reviewers discover limitations that aren't acknowledged in the paper. The authors should use their best judgment and recognize that individual actions in favor of transparency play an important role in developing norms that preserve the integrity of the community. Reviewers will be specifically instructed to not penalize honesty concerning limitations.
    \end{itemize}

\item {\bf Theory assumptions and proofs}
    \item[] Question: For each theoretical result, does the paper provide the full set of assumptions and a complete (and correct) proof?
    \item[] Answer: \answerYes{}
    \item[] Justification: Propositions 1--3 are stated in Section~\ref{sec:method} with all assumptions and notation defined in \S\ref{sec:c1}--\S\ref{sec:c3}, and full proofs are provided in Appendix~\ref{proofs}.
    \item[] Guidelines:
    \begin{itemize}
        \item The answer \answerNA{} means that the paper does not include theoretical results.
        \item All the theorems, formulas, and proofs in the paper should be numbered and cross-referenced.
        \item All assumptions should be clearly stated or referenced in the statement of any theorems.
        \item The proofs can either appear in the main paper or the supplemental material, but if they appear in the supplemental material, the authors are encouraged to provide a short proof sketch to provide intuition.
        \item Inversely, any informal proof provided in the core of the paper should be complemented by formal proofs provided in appendix or supplemental material.
        \item Theorems and Lemmas that the proof relies upon should be properly referenced.
    \end{itemize}

    \item {\bf Experimental result reproducibility}
    \item[] Question: Does the paper fully disclose all the information needed to reproduce the main experimental results of the paper to the extent that it affects the main claims and/or conclusions of the paper (regardless of whether the code and data are provided or not)?
    \item[] Answer: \answerYes{}
    \item[] Justification: Section~\ref{sec:exp} and Appendices~\ref{app:bench}, \ref{app:algo}, \ref{app:baseline} fully document the benchmark construction pipeline, KnowledgeClaw algorithms, per-baseline implementations, and all hyperparameters; an anonymous code release accompanies the submission.
    \item[] Guidelines:
    \begin{itemize}
        \item The answer \answerNA{} means that the paper does not include experiments.
        \item If the paper includes experiments, a \answerNo{} answer to this question will not be perceived well by the reviewers: Making the paper reproducible is important, regardless of whether the code and data are provided or not.
        \item If the contribution is a dataset and\slash or model, the authors should describe the steps taken to make their results reproducible or verifiable.
        \item Depending on the contribution, reproducibility can be accomplished in various ways. For example, if the contribution is a novel architecture, describing the architecture fully might suffice, or if the contribution is a specific model and empirical evaluation, it may be necessary to either make it possible for others to replicate the model with the same dataset, or provide access to the model. In general. releasing code and data is often one good way to accomplish this, but reproducibility can also be provided via detailed instructions for how to replicate the results, access to a hosted model (e.g., in the case of a large language model), releasing of a model checkpoint, or other means that are appropriate to the research performed.
        \item While NeurIPS does not require releasing code, the conference does require all submissions to provide some reasonable avenue for reproducibility, which may depend on the nature of the contribution. For example
        \begin{enumerate}
            \item If the contribution is primarily a new algorithm, the paper should make it clear how to reproduce that algorithm.
            \item If the contribution is primarily a new model architecture, the paper should describe the architecture clearly and fully.
            \item If the contribution is a new model (e.g., a large language model), then there should either be a way to access this model for reproducing the results or a way to reproduce the model (e.g., with an open-source dataset or instructions for how to construct the dataset).
            \item We recognize that reproducibility may be tricky in some cases, in which case authors are welcome to describe the particular way they provide for reproducibility. In the case of closed-source models, it may be that access to the model is limited in some way (e.g., to registered users), but it should be possible for other researchers to have some path to reproducing or verifying the results.
        \end{enumerate}
    \end{itemize}


\item {\bf Open access to data and code}
    \item[] Question: Does the paper provide open access to the data and code, with sufficient instructions to faithfully reproduce the main experimental results, as described in supplemental material?
    \item[] Answer: \answerYes{}
    \item[] Justification: An anonymous GitHub link to the full code, prompts, OmniKE-Bench data, and reproduction scripts is provided in the abstract footnote, with a per-persona SHA-256 manifest documented in Appendix~\ref{app:bench-pipe}.
    \item[] Guidelines:
    \begin{itemize}
        \item The answer \answerNA{} means that paper does not include experiments requiring code.
        \item Please see the NeurIPS code and data submission guidelines (\url{https://neurips.cc/public/guides/CodeSubmissionPolicy}) for more details.
        \item While we encourage the release of code and data, we understand that this might not be possible, so \answerNo{} is an acceptable answer. Papers cannot be rejected simply for not including code, unless this is central to the contribution (e.g., for a new open-source benchmark).
        \item The instructions should contain the exact command and environment needed to run to reproduce the results. See the NeurIPS code and data submission guidelines (\url{https://neurips.cc/public/guides/CodeSubmissionPolicy}) for more details.
        \item The authors should provide instructions on data access and preparation, including how to access the raw data, preprocessed data, intermediate data, and generated data, etc.
        \item The authors should provide scripts to reproduce all experimental results for the new proposed method and baselines. If only a subset of experiments are reproducible, they should state which ones are omitted from the script and why.
        \item At submission time, to preserve anonymity, the authors should release anonymized versions (if applicable).
        \item Providing as much information as possible in supplemental material (appended to the paper) is recommended, but including URLs to data and code is permitted.
    \end{itemize}


\item {\bf Experimental setting/details}
    \item[] Question: Does the paper specify all the training and test details (e.g., data splits, hyperparameters, how they were chosen, type of optimizer) necessary to understand the results?
    \item[] Answer: \answerYes{}
    \item[] Justification: Section~\ref{sec:exp} (Experimental Setup) and Appendix~\ref{app:b-kc} report the unified backbone, embedder, hyperparameters ($\theta_\chi, K_\mathcal{E}, K_\mathcal{H}, \ell, k_0, \gamma, \theta_\pi$), 5-seed protocol, and per-task evaluation rules.
    \item[] Guidelines:
    \begin{itemize}
        \item The answer \answerNA{} means that the paper does not include experiments.
        \item The experimental setting should be presented in the core of the paper to a level of detail that is necessary to appreciate the results and make sense of them.
        \item The full details can be provided either with the code, in appendix, or as supplemental material.
    \end{itemize}

\item {\bf Experiment statistical significance}
    \item[] Question: Does the paper report error bars suitably and correctly defined or other appropriate information about the statistical significance of the experiments?
    \item[] Answer: \answerYes{}
    \item[] Justification: All reported scores are averaged over 5 random seeds per persona (Section~\ref{sec:exp}, Implementation Details), and the lifecycle stress-test in Figure~\ref{fig:cycle} reports mean $\pm$ 1$\sigma$ across these seeds.
    \item[] Guidelines:
    \begin{itemize}
        \item The answer \answerNA{} means that the paper does not include experiments.
        \item The authors should answer \answerYes{} if the results are accompanied by error bars, confidence intervals, or statistical significance tests, at least for the experiments that support the main claims of the paper.
        \item The factors of variability that the error bars are capturing should be clearly stated (for example, train/test split, initialization, random drawing of some parameter, or overall run with given experimental conditions).
        \item The method for calculating the error bars should be explained (closed form formula, call to a library function, bootstrap, etc.)
        \item The assumptions made should be given (e.g., Normally distributed errors).
        \item It should be clear whether the error bar is the standard deviation or the standard error of the mean.
        \item It is OK to report 1-sigma error bars, but one should state it. The authors should preferably report a 2-sigma error bar than state that they have a 96\% CI, if the hypothesis of Normality of errors is not verified.
        \item For asymmetric distributions, the authors should be careful not to show in tables or figures symmetric error bars that would yield results that are out of range (e.g., negative error rates).
        \item If error bars are reported in tables or plots, the authors should explain in the text how they were calculated and reference the corresponding figures or tables in the text.
    \end{itemize}

\item {\bf Experiments compute resources}
    \item[] Question: For each experiment, does the paper provide sufficient information on the computer resources (type of compute workers, memory, time of execution) needed to reproduce the experiments?
    \item[] Answer: \answerYes{}
    \item[] Justification: Section~\ref{sec:exp} and \S\ref{5.7} report the 80-core CPU server, \$1.95 per persona on the OmniKE pipeline, 41-min full-task wall-clock build, and 580\,ms serving latency, with per-method cost--time figures in Figure~\ref{fig:cost}.
    \item[] Guidelines:
    \begin{itemize}
        \item The answer \answerNA{} means that the paper does not include experiments.
        \item The paper should indicate the type of compute workers CPU or GPU, internal cluster, or cloud provider, including relevant memory and storage.
        \item The paper should provide the amount of compute required for each of the individual experimental runs as well as estimate the total compute.
        \item The paper should disclose whether the full research project required more compute than the experiments reported in the paper (e.g., preliminary or failed experiments that didn't make it into the paper).
    \end{itemize}

\item {\bf Code of ethics}
    \item[] Question: Does the research conducted in the paper conform, in every respect, with the NeurIPS Code of Ethics \url{https://neurips.cc/public/EthicsGuidelines}?
    \item[] Answer: \answerYes{}
    \item[] Justification: The research conforms to the NeurIPS Code of Ethics: OmniKE-Bench is constructed from publicly available material with LLM-based annotation, no human subjects are involved, and the submission preserves author and affiliation anonymity.
    \item[] Guidelines:
    \begin{itemize}
        \item The answer \answerNA{} means that the authors have not reviewed the NeurIPS Code of Ethics.
        \item If the authors answer \answerNo, they should explain the special circumstances that require a deviation from the Code of Ethics.
        \item The authors should make sure to preserve anonymity (e.g., if there is a special consideration due to laws or regulations in their jurisdiction).
    \end{itemize}


\item {\bf Broader impacts}
    \item[] Question: Does the paper discuss both potential positive societal impacts and negative societal impacts of the work performed?
    \item[] Answer: \answerYes{}
    \item[] Justification: Appendix~F discusses positive impacts (lowering the cost of live, all-modality personal knowledge engineering) and negative ones (privacy of personal omnimodal corpora and reliance on closed-source LLM/embedding APIs).
    \item[] Guidelines:
    \begin{itemize}
        \item The answer \answerNA{} means that there is no societal impact of the work performed.
        \item If the authors answer \answerNA{} or \answerNo, they should explain why their work has no societal impact or why the paper does not address societal impact.
        \item Examples of negative societal impacts include potential malicious or unintended uses (e.g., disinformation, generating fake profiles, surveillance), fairness considerations (e.g., deployment of technologies that could make decisions that unfairly impact specific groups), privacy considerations, and security considerations.
        \item The conference expects that many papers will be foundational research and not tied to particular applications, let alone deployments. However, if there is a direct path to any negative applications, the authors should point it out. For example, it is legitimate to point out that an improvement in the quality of generative models could be used to generate Deepfakes for disinformation. On the other hand, it is not needed to point out that a generic algorithm for optimizing neural networks could enable people to train models that generate Deepfakes faster.
        \item The authors should consider possible harms that could arise when the technology is being used as intended and functioning correctly, harms that could arise when the technology is being used as intended but gives incorrect results, and harms following from (intentional or unintentional) misuse of the technology.
        \item If there are negative societal impacts, the authors could also discuss possible mitigation strategies (e.g., gated release of models, providing defenses in addition to attacks, mechanisms for monitoring misuse, mechanisms to monitor how a system learns from feedback over time, improving the efficiency and accessibility of ML).
    \end{itemize}

\item {\bf Safeguards}
    \item[] Question: Does the paper describe safeguards that have been put in place for responsible release of data or models that have a high risk for misuse (e.g., pre-trained language models, image generators, or scraped datasets)?
    \item[] Answer: \answerNA{}
    \item[] Justification: The released artefact is a benchmark of synthetic professional personas without sensitive personal data, and KnowledgeClaw is a non-generative agent built on top of public LLM/embedding APIs, posing no high-risk misuse beyond that of the underlying backbones.
    \item[] Guidelines:
    \begin{itemize}
        \item The answer \answerNA{} means that the paper poses no such risks.
        \item Released models that have a high risk for misuse or dual-use should be released with necessary safeguards to allow for controlled use of the model, for example by requiring that users adhere to usage guidelines or restrictions to access the model or implementing safety filters.
        \item Datasets that have been scraped from the Internet could pose safety risks. The authors should describe how they avoided releasing unsafe images.
        \item We recognize that providing effective safeguards is challenging, and many papers do not require this, but we encourage authors to take this into account and make a best faith effort.
    \end{itemize}

\item {\bf Licenses for existing assets}
    \item[] Question: Are the creators or original owners of assets (e.g., code, data, models), used in the paper, properly credited and are the license and terms of use explicitly mentioned and properly respected?
    \item[] Answer: \answerYes{}
    \item[] Justification: All 22 baselines, the Gemini LLM/embedding APIs, and tooling (PyMuPDF, Louvain) are cited and used under their public licenses; per-method citations and version pins are recorded in Appendix~\ref{app:baseline} and Appendix~\ref{app:bench-pipe}.
    \item[] Guidelines:
    \begin{itemize}
        \item The answer \answerNA{} means that the paper does not use existing assets.
        \item The authors should cite the original paper that produced the code package or dataset.
        \item The authors should state which version of the asset is used and, if possible, include a URL.
        \item The name of the license (e.g., CC-BY 4.0) should be included for each asset.
        \item For scraped data from a particular source (e.g., website), the copyright and terms of service of that source should be provided.
        \item If assets are released, the license, copyright information, and terms of use in the package should be provided. For popular datasets, \url{paperswithcode.com/datasets} has curated licenses for some datasets. Their licensing guide can help determine the license of a dataset.
        \item For existing datasets that are re-packaged, both the original license and the license of the derived asset (if it has changed) should be provided.
        \item If this information is not available online, the authors are encouraged to reach out to the asset's creators.
    \end{itemize}

\item {\bf New assets}
    \item[] Question: Are new assets introduced in the paper well documented and is the documentation provided alongside the assets?
    \item[] Answer: \answerYes{}
    \item[] Justification: OmniKE-Bench is released with per-persona corpora, gold annotations, scoring scripts, and a SHA-256 manifest; documentation, task IO signatures, and metric formulas are in Appendix~\ref{app:bench}, and the anonymous repository linked from the abstract footnote.
    \item[] Guidelines:
    \begin{itemize}
        \item The answer \answerNA{} means that the paper does not release new assets.
        \item Researchers should communicate the details of the dataset\slash code\slash model as part of their submissions via structured templates. This includes details about training, license, limitations, etc.
        \item The paper should discuss whether and how consent was obtained from people whose asset is used.
        \item At submission time, remember to anonymize your assets (if applicable). You can either create an anonymized URL or include an anonymized zip file.
    \end{itemize}

\item {\bf Crowdsourcing and research with human subjects}
    \item[] Question: For crowdsourcing experiments and research with human subjects, does the paper include the full text of instructions given to participants and screenshots, if applicable, as well as details about compensation (if any)?
    \item[] Answer: \answerNA{}
    \item[] Justification: The benchmark is constructed from publicly available files with LLM-based annotation (Appendix~\ref{app:bench-pipe}); no crowdsourcing nor human-subject data collection was involved.
    \item[] Guidelines:
    \begin{itemize}
        \item The answer \answerNA{} means that the paper does not involve crowdsourcing nor research with human subjects.
        \item Including this information in the supplemental material is fine, but if the main contribution of the paper involves human subjects, then as much detail as possible should be included in the main paper.
        \item According to the NeurIPS Code of Ethics, workers involved in data collection, curation, or other labor should be paid at least the minimum wage in the country of the data collector.
    \end{itemize}

\item {\bf Institutional review board (IRB) approvals or equivalent for research with human subjects}
    \item[] Question: Does the paper describe potential risks incurred by study participants, whether such risks were disclosed to the subjects, and whether Institutional Review Board (IRB) approvals (or an equivalent approval/review based on the requirements of your country or institution) were obtained?
    \item[] Answer: \answerNA{}
    \item[] Justification: The work involves no human subjects; OmniKE-Bench is constructed from publicly available files with LLM annotation, so IRB review is not applicable.
    \item[] Guidelines:
    \begin{itemize}
        \item The answer \answerNA{} means that the paper does not involve crowdsourcing nor research with human subjects.
        \item Depending on the country in which research is conducted, IRB approval (or equivalent) may be required for any human subjects research. If you obtained IRB approval, you should clearly state this in the paper.
        \item We recognize that the procedures for this may vary significantly between institutions and locations, and we expect authors to adhere to the NeurIPS Code of Ethics and the guidelines for their institution.
        \item For initial submissions, do not include any information that would break anonymity (if applicable), such as the institution conducting the review.
    \end{itemize}

\item {\bf Declaration of LLM usage}
    \item[] Question: Does the paper describe the usage of LLMs if it is an important, original, or non-standard component of the core methods in this research? Note that if the LLM is used only for writing, editing, or formatting purposes and does \emph{not} impact the core methodology, scientific rigor, or originality of the research, declaration is not required.
    %this research?
    \item[] Answer: \answerNA{}
    \item[] Justification: LLMs are used only as a standard backbone for extraction, generation, and evaluation following established RAG/agentic-KE practice, with no original or non-standard LLM component in the core methodology.
    \item[] Guidelines:
    \begin{itemize}
        \item The answer \answerNA{} means that the core method development in this research does not involve LLMs as any important, original, or non-standard components.
        \item Please refer to our LLM policy in the NeurIPS handbook for what should or should not be described.
    \end{itemize}

\end{enumerate}



\end{document}